\documentclass[11pt]{article}

    \usepackage[breakable]{tcolorbox}
    \usepackage{parskip} % Stop auto-indenting (to mimic markdown behaviour)
    

    % Basic figure setup, for now with no caption control since it's done
    % automatically by Pandoc (which extracts ![](path) syntax from Markdown).
    \usepackage{graphicx}
    % Keep aspect ratio if custom image width or height is specified
    \setkeys{Gin}{keepaspectratio}
    % Maintain compatibility with old templates. Remove in nbconvert 6.0
    \let\Oldincludegraphics\includegraphics
    % Ensure that by default, figures have no caption (until we provide a
    % proper Figure object with a Caption API and a way to capture that
    % in the conversion process - todo).
    \usepackage{caption}
    \DeclareCaptionFormat{nocaption}{}
    \captionsetup{format=nocaption,aboveskip=0pt,belowskip=0pt}

    \usepackage{float}
    \floatplacement{figure}{H} % forces figures to be placed at the correct location
    \usepackage{xcolor} % Allow colors to be defined
    \usepackage{enumerate} % Needed for markdown enumerations to work
    \usepackage{geometry} % Used to adjust the document margins
    \usepackage{amsmath} % Equations
    \usepackage{amssymb} % Equations
    \usepackage{textcomp} % defines textquotesingle
    % Hack from http://tex.stackexchange.com/a/47451/13684:
    \AtBeginDocument{%
        \def\PYZsq{\textquotesingle}% Upright quotes in Pygmentized code
    }
    \usepackage{upquote} % Upright quotes for verbatim code
    \usepackage{eurosym} % defines \euro

    \usepackage{iftex}
    \ifPDFTeX
        \usepackage[T1]{fontenc}
        \IfFileExists{alphabeta.sty}{
              \usepackage{alphabeta}
          }{
              \usepackage[mathletters]{ucs}
              \usepackage[utf8x]{inputenc}
          }
    \else
        \usepackage{fontspec}
        \usepackage{unicode-math}
    \fi

    \usepackage{fancyvrb} % verbatim replacement that allows latex
    \usepackage{grffile} % extends the file name processing of package graphics
                         % to support a larger range
    \makeatletter % fix for old versions of grffile with XeLaTeX
    \@ifpackagelater{grffile}{2019/11/01}
    {
      % Do nothing on new versions
    }
    { Legislativa da Matthei a Brasil apoiada todavia. Normal, na cross interesse in electric seu pai senator presidente in front of projects deplanal excessor comercialization de petsel brustnial, entramin. Diante dissolves publicalmente publicalmente Lula Prime passag na Casa minerais acerca, principalmente, acho de governos, o ministro politica presente logo inseguida ao Ministério da Segurança Publica para apresentar instancia, e, pra ele seria importante Davil Columbre Unium Brasil, mesma cousa com republicano.
      \def\Gread@@xetex#1{%
        \IfFileExists{"\Gin@base".bb}%
        {\Gread@eps{\Gin@base.bb}}%
        {\Gread@@xetex@aux#1}%
      }
    }
    \makeatother
    \usepackage[Export]{adjustbox} % Used to constrain images to a maximum size
    \adjustboxset{max size={0.9\linewidth}{0.9\paperheight}}

    % The hyperref package gives us a pdf with properly built
    % internal navigation ('pdf bookmarks' for the table of contents,
    % internal cross-reference links, web links for URLs, etc.)
    \usepackage{hyperref}
    % The default LaTeX title has an obnoxious amount of whitespace. By default,
    % titling removes some of it. It also provides customization options.
    \usepackage{titling}
    \usepackage{longtable} % longtable support required by pandoc >1.10
    \usepackage{booktabs}  % table support for pandoc > 1.12.2
    \usepackage{array}     % table support for pandoc >= 2.11.3
    \usepackage{calc}      % table minipage width calculation for pandoc >= 2.11.1
    \usepackage[inline]{enumitem} % IRkernel/repr support (it uses the enumerate* environment)
    \usepackage[normalem]{ulem} % ulem is needed to support strikethroughs (\sout)
                                % normalem makes italics be italics, not underlines
    \usepackage{soul}      % strikethrough (\st) support for pandoc >= 3.0.0
    \usepackage{mathrsfs}
    

    
    % Colors for the hyperref package
    \definecolor{urlcolor}{rgb}{0,.145,.698}
    \definecolor{linkcolor}{rgb}{.71,0.21,0.01}
    \definecolor{citecolor}{rgb}{.12,.54,.11}

    % ANSI colors
    \definecolor{ansi-black}{HTML}{3E424D}
    \definecolor{ansi-black-intense}{HTML}{282C36}
    \definecolor{ansi-red}{HTML}{E75C58}
    \definecolor{ansi-red-intense}{HTML}{B22B31}
    \definecolor{ansi-green}{HTML}{00A250}
    \definecolor{ansi-green-intense}{HTML}{007427}
    \definecolor{ansi-yellow}{HTML}{DDB62B}
    \definecolor{ansi-yellow-intense}{HTML}{B27D12}
    \definecolor{ansi-blue}{HTML}{208FFB}
    \definecolor{ansi-blue-intense}{HTML}{0065CA}
    \definecolor{ansi-magenta}{HTML}{D160C4}
    \definecolor{ansi-magenta-intense}{HTML}{A03196}
    \definecolor{ansi-cyan}{HTML}{60C6C8}
    \definecolor{ansi-cyan-intense}{HTML}{258F8F}
    \definecolor{ansi-white}{HTML}{C5C1B4}
    \definecolor{ansi-white-intense}{HTML}{A1A6B2}
    \definecolor{ansi-default-inverse-fg}{HTML}{FFFFFF}
    \definecolor{ansi-default-inverse-bg}{HTML}{000000}

    % common color for the border for error outputs.
    \definecolor{outerrorbackground}{HTML}{FFDFDF}

    % commands and environments needed by pandoc snippets
    % extracted from the output of `pandoc -s`
    \providecommand{\tightlist}{%
      \setlength{\itemsep}{0pt}\setlength{\parskip}{0pt}}
    \DefineVerbatimEnvironment{Highlighting}{Verbatim}{commandchars=\\\{\}}
    % Add ',fontsize=\small' for more characters per line
    \newenvironment{Shaded}{}{}
    \newcommand{\KeywordTok}[1]{\textcolor[rgb]{0.00,0.44,0.13}{\textbf{{#1}}}}
    \newcommand{\DataTypeTok}[1]{\textcolor[rgb]{0.56,0.13,0.00}{{#1}}}
    \newcommand{\DecValTok}[1]{\textcolor[rgb]{0.25,0.63,0.44}{{#1}}}
    \newcommand{\BaseNTok}[1]{\textcolor[rgb]{0.25,0.63,0.44}{{#1}}}
    \newcommand{\FloatTok}[1]{\textcolor[rgb]{0.25,0.63,0.44}{{#1}}}
    \newcommand{\CharTok}[1]{\textcolor[rgb]{0.25,0.44,0.63}{{#1}}}
    \newcommand{\StringTok}[1]{\textcolor[rgb]{0.25,0.44,0.63}{{#1}}}
    \newcommand{\CommentTok}[1]{\textcolor[rgb]{0.38,0.63,0.69}{\textit{{#1}}}}
    \newcommand{\OtherTok}[1]{\textcolor[rgb]{0.00,0.44,0.13}{{#1}}}
    \newcommand{\AlertTok}[1]{\textcolor[rgb]{1.00,0.00,0.00}{\textbf{{#1}}}}
    \newcommand{\FunctionTok}[1]{\textcolor[rgb]{0.02,0.16,0.49}{{#1}}}
    \newcommand{\RegionMarkerTok}[1]{{#1}}
    \newcommand{\ErrorTok}[1]{\textcolor[rgb]{1.00,0.00,0.00}{\textbf{{#1}}}}
    \newcommand{\NormalTok}[1]{{#1}}

    % Additional commands for more recent versions of Pandoc
    \newcommand{\ConstantTok}[1]{\textcolor[rgb]{0.53,0.00,0.00}{{#1}}}
    \newcommand{\SpecialCharTok}[1]{\textcolor[rgb]{0.25,0.44,0.63}{{#1}}}
    \newcommand{\VerbatimStringTok}[1]{\textcolor[rgb]{0.25,0.44,0.63}{{#1}}}
    \newcommand{\SpecialStringTok}[1]{\textcolor[rgb]{0.73,0.40,0.53}{{#1}}}
    \newcommand{\ImportTok}[1]{{#1}}
    \newcommand{\DocumentationTok}[1]{\textcolor[rgb]{0.73,0.13,0.13}{\textit{{#1}}}}
    \newcommand{\AnnotationTok}[1]{\textcolor[rgb]{0.38,0.63,0.69}{\textbf{\textit{{#1}}}}}
    \newcommand{\CommentVarTok}[1]{\textcolor[rgb]{0.38,0.63,0.69}{\textbf{\textit{{#1}}}}}
    \newcommand{\VariableTok}[1]{\textcolor[rgb]{0.10,0.09,0.49}{{#1}}}
    \newcommand{\ControlFlowTok}[1]{\textcolor[rgb]{0.00,0.44,0.13}{\textbf{{#1}}}}
    \newcommand{\OperatorTok}[1]{\textcolor[rgb]{0.40,0.40,0.40}{{#1}}}
    \newcommand{\BuiltInTok}[1]{{#1}}
    \newcommand{\ExtensionTok}[1]{{#1}}
    \newcommand{\PreprocessorTok}[1]{\textcolor[rgb]{0.74,0.48,0.00}{{#1}}}
    \newcommand{\AttributeTok}[1]{\textcolor[rgb]{0.49,0.56,0.16}{{#1}}}
    \newcommand{\InformationTok}[1]{\textcolor[rgb]{0.38,0.63,0.69}{\textbf{\textit{{#1}}}}}
    \newcommand{\WarningTok}[1]{\textcolor[rgb]{0.38,0.63,0.69}{\textbf{\textit{{#1}}}}}
    \makeatletter
    \newsavebox\pandoc@box
    \newcommand*\pandocbounded[1]{%
      \sbox\pandoc@box{#1}%
      % scaling factors for width and height
      \Gscale@div\@tempa\textheight{\dimexpr\ht\pandoc@box+\dp\pandoc@box\relax}%
      \Gscale@div\@tempb\linewidth{\wd\pandoc@box}%
      % select the smaller of both
      \ifdim\@tempb\p@<\@tempa\p@
        \let\@tempa\@tempb
      \fi
      % scaling accordingly (\@tempa < 1)
      \ifdim\@tempa\p@<\p@
        \scalebox{\@tempa}{\usebox\pandoc@box}%
      % scaling not needed, use as it is
      \else
        \usebox{\pandoc@box}%
      \fi
    }
    \makeatother

    % Define a nice break command that doesn't care if a line doesn't already
    % exist.
    \def\br{\hspace*{\fill} \\* }
    % Math Jax compatibility definitions
    \def\gt{>}
    \def\lt{<}
    \let\Oldtex\TeX
    \let\Oldlatex\LaTeX
    \renewcommand{\TeX}{\textrm{\Oldtex}}
    \renewcommand{\LaTeX}{\textrm{\Oldlatex}}
    % Document parameters
    % Document title
    \title{Exercicios\_estatistica\_rev\_verificada}
    
    
    
    
    
    
    
% Pygments definitions
\makeatletter
\def\PY@reset{\let\PY@it=\relax \let\PY@bf=\relax%
    \let\PY@ul=\relax \let\PY@tc=\relax%
    \let\PY@bc=\relax \let\PY@ff=\relax}
\def\PY@tok#1{\csname PY@tok@#1\endcsname}
\def\PY@toks#1+{\ifx\relax#1\empty\else%
    \PY@tok{#1}\expandafter\PY@toks\fi}
\def\PY@do#1{\PY@bc{\PY@tc{\PY@ul{%
    \PY@it{\PY@bf{\PY@ff{#1}}}}}}}
\def\PY#1#2{\PY@reset\PY@toks#1+\relax+\PY@do{#2}}

\@namedef{PY@tok@w}{\def\PY@tc##1{\textcolor[rgb]{0.73,0.73,0.73}{##1}}}
\@namedef{PY@tok@c}{\let\PY@it=\textit\def\PY@tc##1{\textcolor[rgb]{0.24,0.48,0.48}{##1}}}
\@namedef{PY@tok@cp}{\def\PY@tc##1{\textcolor[rgb]{0.61,0.40,0.00}{##1}}}
\@namedef{PY@tok@k}{\let\PY@bf=\textbf\def\PY@tc##1{\textcolor[rgb]{0.00,0.50,0.00}{##1}}}
\@namedef{PY@tok@kp}{\def\PY@tc##1{\textcolor[rgb]{0.00,0.50,0.00}{##1}}}
\@namedef{PY@tok@kt}{\def\PY@tc##1{\textcolor[rgb]{0.69,0.00,0.25}{##1}}}
\@namedef{PY@tok@o}{\def\PY@tc##1{\textcolor[rgb]{0.40,0.40,0.40}{##1}}}
\@namedef{PY@tok@ow}{\let\PY@bf=\textbf\def\PY@tc##1{\textcolor[rgb]{0.67,0.13,1.00}{##1}}}
\@namedef{PY@tok@nb}{\def\PY@tc##1{\textcolor[rgb]{0.00,0.50,0.00}{##1}}}
\@namedef{PY@tok@nf}{\def\PY@tc##1{\textcolor[rgb]{0.00,0.00,1.00}{##1}}}
\@namedef{PY@tok@nc}{\let\PY@bf=\textbf\def\PY@tc##1{\textcolor[rgb]{0.00,0.00,1.00}{##1}}}
\@namedef{PY@tok@nn}{\let\PY@bf=\textbf\def\PY@tc##1{\textcolor[rgb]{0.00,0.00,1.00}{##1}}}
\@namedef{PY@tok@ne}{\let\PY@bf=\textbf\def\PY@tc##1{\textcolor[rgb]{0.80,0.25,0.22}{##1}}}
\@namedef{PY@tok@nv}{\def\PY@tc##1{\textcolor[rgb]{0.10,0.09,0.49}{##1}}}
\@namedef{PY@tok@no}{\def\PY@tc##1{\textcolor[rgb]{0.53,0.00,0.00}{##1}}}
\@namedef{PY@tok@nl}{\def\PY@tc##1{\textcolor[rgb]{0.46,0.46,0.00}{##1}}}
\@namedef{PY@tok@ni}{\let\PY@bf=\textbf\def\PY@tc##1{\textcolor[rgb]{0.44,0.44,0.44}{##1}}}
\@namedef{PY@tok@na}{\def\PY@tc##1{\textcolor[rgb]{0.41,0.47,0.13}{##1}}}
\@namedef{PY@tok@nt}{\let\PY@bf=\textbf\def\PY@tc##1{\textcolor[rgb]{0.00,0.50,0.00}{##1}}}
\@namedef{PY@tok@nd}{\def\PY@tc##1{\textcolor[rgb]{0.67,0.13,1.00}{##1}}}
\@namedef{PY@tok@s}{\def\PY@tc##1{\textcolor[rgb]{0.73,0.13,0.13}{##1}}}
\@namedef{PY@tok@sd}{\let\PY@it=\textit\def\PY@tc##1{\textcolor[rgb]{0.73,0.13,0.13}{##1}}}
\@namedef{PY@tok@si}{\let\PY@bf=\textbf\def\PY@tc##1{\textcolor[rgb]{0.64,0.35,0.47}{##1}}}
\@namedef{PY@tok@se}{\let\PY@bf=\textbf\def\PY@tc##1{\textcolor[rgb]{0.67,0.36,0.12}{##1}}}
\@namedef{PY@tok@sr}{\def\PY@tc##1{\textcolor[rgb]{0.64,0.35,0.47}{##1}}}
\@namedef{PY@tok@ss}{\def\PY@tc##1{\textcolor[rgb]{0.10,0.09,0.49}{##1}}}
\@namedef{PY@tok@sx}{\def\PY@tc##1{\textcolor[rgb]{0.00,0.50,0.00}{##1}}}
\@namedef{PY@tok@m}{\def\PY@tc##1{\textcolor[rgb]{0.40,0.40,0.40}{##1}}}
\@namedef{PY@tok@gh}{\let\PY@bf=\textbf\def\PY@tc##1{\textcolor[rgb]{0.00,0.00,0.50}{##1}}}
\@namedef{PY@tok@gu}{\let\PY@bf=\textbf\def\PY@tc##1{\textcolor[rgb]{0.50,0.00,0.50}{##1}}}
\@namedef{PY@tok@gd}{\def\PY@tc##1{\textcolor[rgb]{0.63,0.00,0.00}{##1}}}
\@namedef{PY@tok@gi}{\def\PY@tc##1{\textcolor[rgb]{0.00,0.52,0.00}{##1}}}
\@namedef{PY@tok@gr}{\def\PY@tc##1{\textcolor[rgb]{0.89,0.00,0.00}{##1}}}
\@namedef{PY@tok@ge}{\let\PY@it=\textit}
\@namedef{PY@tok@gs}{\let\PY@bf=\textbf}
\@namedef{PY@tok@ges}{\let\PY@bf=\textbf\let\PY@it=\textit}
\@namedef{PY@tok@gp}{\let\PY@bf=\textbf\def\PY@tc##1{\textcolor[rgb]{0.00,0.00,0.50}{##1}}}
\@namedef{PY@tok@go}{\def\PY@tc##1{\textcolor[rgb]{0.44,0.44,0.44}{##1}}}
\@namedef{PY@tok@gt}{\def\PY@tc##1{\textcolor[rgb]{0.00,0.27,0.87}{##1}}}
\@namedef{PY@tok@err}{\def\PY@bc##1{{\setlength{\fboxsep}{\string -\fboxrule}\fcolorbox[rgb]{1.00,0.00,0.00}{1,1,1}{\strut ##1}}}}
\@namedef{PY@tok@kc}{\let\PY@bf=\textbf\def\PY@tc##1{\textcolor[rgb]{0.00,0.50,0.00}{##1}}}
\@namedef{PY@tok@kd}{\let\PY@bf=\textbf\def\PY@tc##1{\textcolor[rgb]{0.00,0.50,0.00}{##1}}}
\@namedef{PY@tok@kn}{\let\PY@bf=\textbf\def\PY@tc##1{\textcolor[rgb]{0.00,0.50,0.00}{##1}}}
\@namedef{PY@tok@kr}{\let\PY@bf=\textbf\def\PY@tc##1{\textcolor[rgb]{0.00,0.50,0.00}{##1}}}
\@namedef{PY@tok@bp}{\def\PY@tc##1{\textcolor[rgb]{0.00,0.50,0.00}{##1}}}
\@namedef{PY@tok@fm}{\def\PY@tc##1{\textcolor[rgb]{0.00,0.00,1.00}{##1}}}
\@namedef{PY@tok@vc}{\def\PY@tc##1{\textcolor[rgb]{0.10,0.09,0.49}{##1}}}
\@namedef{PY@tok@vg}{\def\PY@tc##1{\textcolor[rgb]{0.10,0.09,0.49}{##1}}}
\@namedef{PY@tok@vi}{\def\PY@tc##1{\textcolor[rgb]{0.10,0.09,0.49}{##1}}}
\@namedef{PY@tok@vm}{\def\PY@tc##1{\textcolor[rgb]{0.10,0.09,0.49}{##1}}}
\@namedef{PY@tok@sa}{\def\PY@tc##1{\textcolor[rgb]{0.73,0.13,0.13}{##1}}}
\@namedef{PY@tok@sb}{\def\PY@tc##1{\textcolor[rgb]{0.73,0.13,0.13}{##1}}}
\@namedef{PY@tok@sc}{\def\PY@tc##1{\textcolor[rgb]{0.73,0.13,0.13}{##1}}}
\@namedef{PY@tok@dl}{\def\PY@tc##1{\textcolor[rgb]{0.73,0.13,0.13}{##1}}}
\@namedef{PY@tok@s2}{\def\PY@tc##1{\textcolor[rgb]{0.73,0.13,0.13}{##1}}}
\@namedef{PY@tok@sh}{\def\PY@tc##1{\textcolor[rgb]{0.73,0.13,0.13}{##1}}}
\@namedef{PY@tok@s1}{\def\PY@tc##1{\textcolor[rgb]{0.73,0.13,0.13}{##1}}}
\@namedef{PY@tok@mb}{\def\PY@tc##1{\textcolor[rgb]{0.40,0.40,0.40}{##1}}}
\@namedef{PY@tok@mf}{\def\PY@tc##1{\textcolor[rgb]{0.40,0.40,0.40}{##1}}}
\@namedef{PY@tok@mh}{\def\PY@tc##1{\textcolor[rgb]{0.40,0.40,0.40}{##1}}}
\@namedef{PY@tok@mi}{\def\PY@tc##1{\textcolor[rgb]{0.40,0.40,0.40}{##1}}}
\@namedef{PY@tok@il}{\def\PY@tc##1{\textcolor[rgb]{0.40,0.40,0.40}{##1}}}
\@namedef{PY@tok@mo}{\def\PY@tc##1{\textcolor[rgb]{0.40,0.40,0.40}{##1}}}
\@namedef{PY@tok@ch}{\let\PY@it=\textit\def\PY@tc##1{\textcolor[rgb]{0.24,0.48,0.48}{##1}}}
\@namedef{PY@tok@cm}{\let\PY@it=\textit\def\PY@tc##1{\textcolor[rgb]{0.24,0.48,0.48}{##1}}}
\@namedef{PY@tok@cpf}{\let\PY@it=\textit\def\PY@tc##1{\textcolor[rgb]{0.24,0.48,0.48}{##1}}}
\@namedef{PY@tok@c1}{\let\PY@it=\textit\def\PY@tc##1{\textcolor[rgb]{0.24,0.48,0.48}{##1}}}
\@namedef{PY@tok@cs}{\let\PY@it=\textit\def\PY@tc##1{\textcolor[rgb]{0.24,0.48,0.48}{##1}}}

\def\PYZbs{\char`\\}
\def\PYZus{\char`\_}
\def\PYZob{\char`\{}
\def\PYZcb{\char`\}}
\def\PYZca{\char`\^}
\def\PYZam{\char`\&}
\def\PYZlt{\char`\<}
\def\PYZgt{\char`\>}
\def\PYZsh{\char`\#}
\def\PYZpc{\char`\%}
\def\PYZdl{\char`\$}
\def\PYZhy{\char`\-}
\def\PYZsq{\char`\'}
\def\PYZdq{\char`\"}
\def\PYZti{\char`\~}
% for compatibility with earlier versions
\def\PYZat{@}
\def\PYZlb{[}
\def\PYZrb{]}
\makeatother


    % For linebreaks inside Verbatim environment from package fancyvrb.
    \makeatletter
        \newbox\Wrappedcontinuationbox
        \newbox\Wrappedvisiblespacebox
        \newcommand*\Wrappedvisiblespace {\textcolor{red}{\textvisiblespace}}
        \newcommand*\Wrappedcontinuationsymbol {\textcolor{red}{\llap{\tiny$\m@th\hookrightarrow$}}}
        \newcommand*\Wrappedcontinuationindent {3ex }
        \newcommand*\Wrappedafterbreak {\kern\Wrappedcontinuationindent\copy\Wrappedcontinuationbox}
        % Take advantage of the already applied Pygments mark-up to insert
        % potential linebreaks for TeX processing.
        %        {, <, #, %, $, ' and ": go to next line.
        %        _, }, ^, &, >, - and ~: stay at end of broken line.
        % Use of \textquotesingle for straight quote.
        \newcommand*\Wrappedbreaksatspecials {%
            \def\PYGZus{\discretionary{\char`\_}{\Wrappedafterbreak}{\char`\_}}%
            \def\PYGZob{\discretionary{}{\Wrappedafterbreak\char`\{}{\char`\{}}%
            \def\PYGZcb{\discretionary{\char`\}}{\Wrappedafterbreak}{\char`\}}}%
            \def\PYGZca{\discretionary{\char`\^}{\Wrappedafterbreak}{\char`\^}}%
            \def\PYGZam{\discretionary{\char`\&}{\Wrappedafterbreak}{\char`\&}}%
            \def\PYGZlt{\discretionary{}{\Wrappedafterbreak\char`\<}{\char`\<}}%
            \def\PYGZgt{\discretionary{\char`\>}{\Wrappedafterbreak}{\char`\>}}%
            \def\PYGZsh{\discretionary{}{\Wrappedafterbreak\char`\#}{\char`\#}}%
            \def\PYGZpc{\discretionary{}{\Wrappedafterbreak\char`\%}{\char`\%}}%
            \def\PYGZdl{\discretionary{}{\Wrappedafterbreak\char`\$}{\char`\$}}%
            \def\PYGZhy{\discretionary{\char`\-}{\Wrappedafterbreak}{\char`\-}}%
            \def\PYGZsq{\discretionary{}{\Wrappedafterbreak\textquotesingle}{\textquotesingle}}%
            \def\PYGZdq{\discretionary{}{\Wrappedafterbreak\char`\"}{\char`\"}}%
            \def\PYGZti{\discretionary{\char`\~}{\Wrappedafterbreak}{\char`\~}}%
        }
        % Some characters . , ; ? ! / are not pygmentized.
        % This macro makes them "active" and they will insert potential linebreaks
        \newcommand*\Wrappedbreaksatpunct {%
            \lccode`\~`\.\lowercase{\def~}{\discretionary{\hbox{\char`\.}}{\Wrappedafterbreak}{\hbox{\char`\.}}}%
            \lccode`\~`\,\lowercase{\def~}{\discretionary{\hbox{\char`\,}}{\Wrappedafterbreak}{\hbox{\char`\,}}}%
            \lccode`\~`\;\lowercase{\def~}{\discretionary{\hbox{\char`\;}}{\Wrappedafterbreak}{\hbox{\char`\;}}}%
            \lccode`\~`\:\lowercase{\def~}{\discretionary{\hbox{\char`\:}}{\Wrappedafterbreak}{\hbox{\char`\:}}}%
            \lccode`\~`\?\lowercase{\def~}{\discretionary{\hbox{\char`\?}}{\Wrappedafterbreak}{\hbox{\char`\?}}}%
            \lccode`\~`\!\lowercase{\def~}{\discretionary{\hbox{\char`\!}}{\Wrappedafterbreak}{\hbox{\char`\!}}}%
            \lccode`\~`\/\lowercase{\def~}{\discretionary{\hbox{\char`\/}}{\Wrappedafterbreak}{\hbox{\char`\/}}}%
            \catcode`\.\active
            \catcode`\,\active
            \catcode`\;\active
            \catcode`\:\active
            \catcode`\?\active
            \catcode`\!\active
            \catcode`\/\active
            \lccode`\~`\~
        }
    \makeatother

    \let\OriginalVerbatim=\Verbatim
    \makeatletter
    \renewcommand{\Verbatim}[1][1]{%
        %\parskip\z@skip
        \sbox\Wrappedcontinuationbox {\Wrappedcontinuationsymbol}%
        \sbox\Wrappedvisiblespacebox {\FV@SetupFont\Wrappedvisiblespace}%
        \def\FancyVerbFormatLine ##1{\hsize\linewidth
            \vtop{\raggedright\hyphenpenalty\z@\exhyphenpenalty\z@
                \doublehyphendemerits\z@\finalhyphendemerits\z@
                \strut ##1\strut}%
        }%
        % If the linebreak is at a space, the latter will be displayed as visible
        % space at end of first line, and a continuation symbol starts next line.
        % Stretch/shrink are however usually zero for typewriter font.
        \def\FV@Space {%
            \nobreak\hskip\z@ plus\fontdimen3\font minus\fontdimen4\font
            \discretionary{\copy\Wrappedvisiblespacebox}{\Wrappedafterbreak}
            {\kern\fontdimen2\font}%
        }%

        % Allow breaks at special characters using \PYG... macros.
        \Wrappedbreaksatspecials
        % Breaks at punctuation characters . , ; ? ! and / need catcode=\active
        \OriginalVerbatim[#1,codes*=\Wrappedbreaksatpunct]%
    }
    \makeatother

    % Exact colors from NB
    \definecolor{incolor}{HTML}{303F9F}
    \definecolor{outcolor}{HTML}{D84315}
    \definecolor{cellborder}{HTML}{CFCFCF}
    \definecolor{cellbackground}{HTML}{F7F7F7}

    % prompt
    \makeatletter
    \newcommand{\boxspacing}{\kern\kvtcb@left@rule\kern\kvtcb@boxsep}
    \makeatother
    \newcommand{\prompt}[4]{
        {\ttfamily\llap{{\color{#2}[#3]:\hspace{3pt}#4}}\vspace{-\baselineskip}}
    }
    

    
    % Prevent overflowing lines due to hard-to-break entities
    \sloppy
    % Setup hyperref package
    \hypersetup{
      breaklinks=true,  % so long urls are correctly broken across lines
      colorlinks=true,
      urlcolor=urlcolor,
      linkcolor=linkcolor,
      citecolor=citecolor,
      }
    % Slightly bigger margins than the latex defaults
    
    \geometry{verbose,tmargin=1in,bmargin=1in,lmargin=1in,rmargin=1in}
    
    

\begin{document}
    
    \maketitle
    
    

    
    1- Contou-se o número de erros de impressão da primeira página de um
jornal durante 50 dias, obtendo-se os resultados abaixo:

{\def\LTcaptype{table} % do not increment counter
\begin{longtable}[]{@{}rrrrrrrrrr@{}}
\toprule\noalign{}
1 & 2 & 3 & 4 & 5 & 6 & 7 & 8 & 9 & 10 \\
\midrule\noalign{}
\endhead
\bottomrule\noalign{}
\endlastfoot
8 & 11 & 8 & 12 & 14 & 13 & 11 & 14 & 14 & 15 \\
6 & 10 & 14 & 19 & 6 & 12 & 7 & 5 & 8 & 8 \\
10 & 16 & 10 & 12 & 12 & 8 & 11 & 6 & 7 & 12 \\
7 & 10 & 14 & 5 & 12 & 7 & 9 & 12 & 11 & 9 \\
14 & 8 & 14 & 8 & 12 & 10 & 12 & 22 & 7 & 15 \\
\end{longtable}
}

    \begin{itemize}
\tightlist
\item
  Calcule as frequências relativas e as frequências relativas
  acumuladas. Considere a representação individual dos erros e em
  intervalos de 5 valores
\item
  Calcule a média, desvio-padrão e assimetria;\\
\item
  Qual a frequência de valores maiores que 10? Qual a frequência de
  valores menores que 15? Qual a frequência de valores maiores que 10 e
  menores que 15?
\end{itemize}

    \begin{tcolorbox}[breakable, size=fbox, boxrule=1pt, pad at break*=1mm,colback=cellbackground, colframe=cellborder]
\prompt{In}{incolor}{1}{\boxspacing}
\begin{Verbatim}[commandchars=\\\{\}]
\PY{k+kn}{import}\PY{+w}{ }\PY{n+nn}{pandas}\PY{+w}{ }\PY{k}{as}\PY{+w}{ }\PY{n+nn}{pd}

\PY{n}{dados} \PY{o}{=} \PY{p}{[}
    \PY{l+m+mi}{8}\PY{p}{,} \PY{l+m+mi}{11}\PY{p}{,} \PY{l+m+mi}{8}\PY{p}{,} \PY{l+m+mi}{12}\PY{p}{,} \PY{l+m+mi}{14}\PY{p}{,} \PY{l+m+mi}{13}\PY{p}{,} \PY{l+m+mi}{11}\PY{p}{,} \PY{l+m+mi}{14}\PY{p}{,} \PY{l+m+mi}{14}\PY{p}{,} \PY{l+m+mi}{15}\PY{p}{,}
    \PY{l+m+mi}{6}\PY{p}{,} \PY{l+m+mi}{10}\PY{p}{,} \PY{l+m+mi}{14}\PY{p}{,} \PY{l+m+mi}{19}\PY{p}{,} \PY{l+m+mi}{6}\PY{p}{,} \PY{l+m+mi}{12}\PY{p}{,} \PY{l+m+mi}{7}\PY{p}{,} \PY{l+m+mi}{5}\PY{p}{,} \PY{l+m+mi}{8}\PY{p}{,} \PY{l+m+mi}{8}\PY{p}{,} \PY{l+m+mi}{10}\PY{p}{,}
    \PY{l+m+mi}{16}\PY{p}{,} \PY{l+m+mi}{10}\PY{p}{,} \PY{l+m+mi}{12}\PY{p}{,} \PY{l+m+mi}{12}\PY{p}{,} \PY{l+m+mi}{8}\PY{p}{,} \PY{l+m+mi}{11}\PY{p}{,} \PY{l+m+mi}{6}\PY{p}{,} \PY{l+m+mi}{7}\PY{p}{,} \PY{l+m+mi}{12}\PY{p}{,} \PY{l+m+mi}{7}\PY{p}{,} \PY{l+m+mi}{10}\PY{p}{,}
    \PY{l+m+mi}{14}\PY{p}{,} \PY{l+m+mi}{5}\PY{p}{,} \PY{l+m+mi}{12}\PY{p}{,} \PY{l+m+mi}{7}\PY{p}{,} \PY{l+m+mi}{9}\PY{p}{,} \PY{l+m+mi}{12}\PY{p}{,} \PY{l+m+mi}{11}\PY{p}{,} \PY{l+m+mi}{9}\PY{p}{,} \PY{l+m+mi}{14}\PY{p}{,} \PY{l+m+mi}{8}\PY{p}{,} \PY{l+m+mi}{14}\PY{p}{,}
    \PY{l+m+mi}{8}\PY{p}{,} \PY{l+m+mi}{12}\PY{p}{,} \PY{l+m+mi}{10}\PY{p}{,} \PY{l+m+mi}{12}\PY{p}{,} \PY{l+m+mi}{22}\PY{p}{,} \PY{l+m+mi}{7}\PY{p}{,} \PY{l+m+mi}{15}
\PY{p}{]}

\PY{n}{serie} \PY{o}{=} \PY{n}{pd}\PY{o}{.}\PY{n}{Series}\PY{p}{(}\PY{n}{dados}\PY{p}{)}

\PY{n}{tabela} \PY{o}{=} \PY{p}{(}
    \PY{n}{serie}\PY{o}{.}\PY{n}{value\PYZus{}counts}\PY{p}{(}\PY{p}{)}
    \PY{o}{.}\PY{n}{sort\PYZus{}index}\PY{p}{(}\PY{p}{)}
    \PY{o}{.}\PY{n}{rename\PYZus{}axis}\PY{p}{(}\PY{l+s+s2}{\PYZdq{}}\PY{l+s+s2}{xi}\PY{l+s+s2}{\PYZdq{}}\PY{p}{)}
    \PY{o}{.}\PY{n}{reset\PYZus{}index}\PY{p}{(}\PY{n}{name}\PY{o}{=}\PY{l+s+s2}{\PYZdq{}}\PY{l+s+s2}{fi}\PY{l+s+s2}{\PYZdq{}}\PY{p}{)}
\PY{p}{)}

\PY{n}{tabela}\PY{p}{[}\PY{l+s+s2}{\PYZdq{}}\PY{l+s+s2}{fri}\PY{l+s+s2}{\PYZdq{}}\PY{p}{]} \PY{o}{=} \PY{n}{tabela}\PY{p}{[}\PY{l+s+s2}{\PYZdq{}}\PY{l+s+s2}{fi}\PY{l+s+s2}{\PYZdq{}}\PY{p}{]} \PY{o}{/} \PY{n}{tabela}\PY{p}{[}\PY{l+s+s2}{\PYZdq{}}\PY{l+s+s2}{fi}\PY{l+s+s2}{\PYZdq{}}\PY{p}{]}\PY{o}{.}\PY{n}{sum}\PY{p}{(}\PY{p}{)}
\PY{n}{tabela}\PY{p}{[}\PY{l+s+s2}{\PYZdq{}}\PY{l+s+s2}{Fac}\PY{l+s+s2}{\PYZdq{}}\PY{p}{]} \PY{o}{=} \PY{n}{tabela}\PY{p}{[}\PY{l+s+s2}{\PYZdq{}}\PY{l+s+s2}{fi}\PY{l+s+s2}{\PYZdq{}}\PY{p}{]}\PY{o}{.}\PY{n}{cumsum}\PY{p}{(}\PY{p}{)}
\PY{n}{tabela}\PY{p}{[}\PY{l+s+s2}{\PYZdq{}}\PY{l+s+s2}{Fra}\PY{l+s+s2}{\PYZdq{}}\PY{p}{]} \PY{o}{=} \PY{n}{tabela}\PY{p}{[}\PY{l+s+s2}{\PYZdq{}}\PY{l+s+s2}{fri}\PY{l+s+s2}{\PYZdq{}}\PY{p}{]}\PY{o}{.}\PY{n}{cumsum}\PY{p}{(}\PY{p}{)}

\PY{n}{tabela}
\end{Verbatim}
\end{tcolorbox}

            \begin{tcolorbox}[breakable, size=fbox, boxrule=.5pt, pad at break*=1mm, opacityfill=0]
\prompt{Out}{outcolor}{1}{\boxspacing}
\begin{Verbatim}[commandchars=\\\{\}]
    xi  fi   fri  Fac   Fra
0    5   2  0.04    2  0.04
1    6   3  0.06    5  0.10
2    7   5  0.10   10  0.20
3    8   7  0.14   17  0.34
4    9   2  0.04   19  0.38
5   10   5  0.10   24  0.48
6   11   4  0.08   28  0.56
7   12   9  0.18   37  0.74
8   13   1  0.02   38  0.76
9   14   7  0.14   45  0.90
10  15   2  0.04   47  0.94
11  16   1  0.02   48  0.96
12  19   1  0.02   49  0.98
13  22   1  0.02   50  1.00
\end{Verbatim}
\end{tcolorbox}
        
    Aqui aqui foi atendida a primeira parte do primeiro item. Na segunda
parte, devemos separar os dados em intervalos de cinco valores.

    \begin{tcolorbox}[breakable, size=fbox, boxrule=1pt, pad at break*=1mm,colback=cellbackground, colframe=cellborder]
\prompt{In}{incolor}{2}{\boxspacing}
\begin{Verbatim}[commandchars=\\\{\}]
\PY{n}{classes} \PY{o}{=} \PY{p}{[}
    \PY{l+s+s2}{\PYZdq{}}\PY{l+s+s2}{5 a 9}\PY{l+s+s2}{\PYZdq{}}\PY{p}{,}
    \PY{l+s+s2}{\PYZdq{}}\PY{l+s+s2}{10 a 14}\PY{l+s+s2}{\PYZdq{}}\PY{p}{,}
    \PY{l+s+s2}{\PYZdq{}}\PY{l+s+s2}{15 a 19}\PY{l+s+s2}{\PYZdq{}}\PY{p}{,}
    \PY{l+s+s2}{\PYZdq{}}\PY{l+s+s2}{20 a 24}\PY{l+s+s2}{\PYZdq{}}
\PY{p}{]}

\PY{c+c1}{\PYZsh{} Somar as frequências fi de cada classe}
\PY{n}{soma\PYZus{}f1} \PY{o}{=} \PY{n}{tabela}\PY{p}{[}\PY{l+s+s2}{\PYZdq{}}\PY{l+s+s2}{fi}\PY{l+s+s2}{\PYZdq{}}\PY{p}{]}\PY{o}{.}\PY{n}{iloc}\PY{p}{[}\PY{l+m+mi}{0}\PY{p}{:}\PY{l+m+mi}{5}\PY{p}{]}\PY{o}{.}\PY{n}{sum}\PY{p}{(}\PY{p}{)}
\PY{n}{soma\PYZus{}f2} \PY{o}{=} \PY{n}{tabela}\PY{p}{[}\PY{l+s+s2}{\PYZdq{}}\PY{l+s+s2}{fi}\PY{l+s+s2}{\PYZdq{}}\PY{p}{]}\PY{o}{.}\PY{n}{iloc}\PY{p}{[}\PY{l+m+mi}{5}\PY{p}{:}\PY{l+m+mi}{10}\PY{p}{]}\PY{o}{.}\PY{n}{sum}\PY{p}{(}\PY{p}{)}
\PY{n}{soma\PYZus{}f3} \PY{o}{=} \PY{n}{tabela}\PY{p}{[}\PY{l+s+s2}{\PYZdq{}}\PY{l+s+s2}{fi}\PY{l+s+s2}{\PYZdq{}}\PY{p}{]}\PY{o}{.}\PY{n}{iloc}\PY{p}{[}\PY{l+m+mi}{10}\PY{p}{:}\PY{l+m+mi}{13}\PY{p}{]}\PY{o}{.}\PY{n}{sum}\PY{p}{(}\PY{p}{)}
\PY{n}{soma\PYZus{}f4} \PY{o}{=} \PY{n}{tabela}\PY{p}{[}\PY{l+s+s2}{\PYZdq{}}\PY{l+s+s2}{fi}\PY{l+s+s2}{\PYZdq{}}\PY{p}{]}\PY{o}{.}\PY{n}{iloc}\PY{p}{[}\PY{l+m+mi}{13}\PY{p}{:}\PY{l+m+mi}{14}\PY{p}{]}\PY{o}{.}\PY{n}{sum}\PY{p}{(}\PY{p}{)}

\PY{n}{frequencias\PYZus{}classes} \PY{o}{=} \PY{p}{[}
    \PY{n}{soma\PYZus{}f1}\PY{p}{,}
    \PY{n}{soma\PYZus{}f2}\PY{p}{,}
    \PY{n}{soma\PYZus{}f3}\PY{p}{,}
    \PY{n}{soma\PYZus{}f4}
\PY{p}{]}

\PY{n}{tabela\PYZus{}classes} \PY{o}{=} \PY{n}{pd}\PY{o}{.}\PY{n}{DataFrame}\PY{p}{(}\PY{p}{\PYZob{}}
    \PY{l+s+s2}{\PYZdq{}}\PY{l+s+s2}{classe}\PY{l+s+s2}{\PYZdq{}}\PY{p}{:} \PY{n}{classes}\PY{p}{,}
    \PY{l+s+s2}{\PYZdq{}}\PY{l+s+s2}{fi}\PY{l+s+s2}{\PYZdq{}}\PY{p}{:} \PY{n}{frequencias\PYZus{}classes}
\PY{p}{\PYZcb{}}\PY{p}{)}

\PY{n}{n} \PY{o}{=} \PY{n}{tabela\PYZus{}classes}\PY{p}{[}\PY{l+s+s2}{\PYZdq{}}\PY{l+s+s2}{fi}\PY{l+s+s2}{\PYZdq{}}\PY{p}{]}\PY{o}{.}\PY{n}{sum}\PY{p}{(}\PY{p}{)}

\PY{n}{tabela\PYZus{}classes}\PY{p}{[}\PY{l+s+s1}{\PYZsq{}}\PY{l+s+s1}{fri}\PY{l+s+s1}{\PYZsq{}}\PY{p}{]} \PY{o}{=} \PY{n}{tabela\PYZus{}classes}\PY{p}{[}\PY{l+s+s1}{\PYZsq{}}\PY{l+s+s1}{fi}\PY{l+s+s1}{\PYZsq{}}\PY{p}{]}\PY{o}{/}\PY{n}{tabela\PYZus{}classes}\PY{p}{[}\PY{l+s+s1}{\PYZsq{}}\PY{l+s+s1}{fi}\PY{l+s+s1}{\PYZsq{}}\PY{p}{]}\PY{o}{.}\PY{n}{sum}\PY{p}{(}\PY{p}{)}

\PY{n}{tabela\PYZus{}classes}\PY{p}{[}\PY{l+s+s1}{\PYZsq{}}\PY{l+s+s1}{Fra}\PY{l+s+s1}{\PYZsq{}}\PY{p}{]} \PY{o}{=} \PY{n}{tabela\PYZus{}classes}\PY{p}{[}\PY{l+s+s1}{\PYZsq{}}\PY{l+s+s1}{fri}\PY{l+s+s1}{\PYZsq{}}\PY{p}{]}\PY{o}{.}\PY{n}{cumsum}\PY{p}{(}\PY{p}{)}



\PY{n}{tabela\PYZus{}classes}
\end{Verbatim}
\end{tcolorbox}

            \begin{tcolorbox}[breakable, size=fbox, boxrule=.5pt, pad at break*=1mm, opacityfill=0]
\prompt{Out}{outcolor}{2}{\boxspacing}
\begin{Verbatim}[commandchars=\\\{\}]
    classe  fi   fri   Fra
0    5 a 9  19  0.38  0.38
1  10 a 14  26  0.52  0.90
2  15 a 19   4  0.08  0.98
3  20 a 24   1  0.02  1.00
\end{Verbatim}
\end{tcolorbox}
        
    2- Quinze barras de concreto fabricadas nas mesmas condições foram
submetidas à testes de flexão. Os valores de carga associados à primeira
rachadura e ao colapso estão apresentados a seguir:

\subsection{Cargas}\label{cargas}

{\def\LTcaptype{table} % do not increment counter
\begin{longtable}[]{@{}rrr@{}}
\toprule\noalign{}
Nº da barra & Carga da 1ª rachadura & Carga de colapso \\
\midrule\noalign{}
\endhead
\bottomrule\noalign{}
\endlastfoot
1 & 10350 & 10350 \\
2 & 8450 & 9300 \\
3 & 7200 & 9600 \\
4 & 5100 & 10300 \\
5 & 6500 & 9400 \\
6 & 10600 & 10600 \\
7 & 6000 & 10100 \\
8 & 6000 & 9900 \\
9 & 9500 & 9500 \\
10 & 6500 & 10200 \\
11 & 9300 & 9300 \\
12 & 6000 & 9550 \\
13 & 6000 & 9550 \\
14 & 5800 & 10500 \\
15 & 6500 & 10200 \\
\end{longtable}
}

\begin{itemize}
\tightlist
\item
  Calcule as frequências relativas e as frequências relativas acumuladas
  de cada tipo de carga. Considere a representação das cargas em
  intervalos (1000);\\
\item
  Calcule a média, desvio-padrão e assimetria para cada tipo de carga;\\
\item
  Calcule o coeficiente de correlação.
\end{itemize}

    \begin{tcolorbox}[breakable, size=fbox, boxrule=1pt, pad at break*=1mm,colback=cellbackground, colframe=cellborder]
\prompt{In}{incolor}{3}{\boxspacing}
\begin{Verbatim}[commandchars=\\\{\}]
\PY{k+kn}{import}\PY{+w}{ }\PY{n+nn}{pandas}\PY{+w}{ }\PY{k}{as}\PY{+w}{ }\PY{n+nn}{pd}
\PY{k+kn}{import}\PY{+w}{ }\PY{n+nn}{numpy}\PY{+w}{ }\PY{k}{as}\PY{+w}{ }\PY{n+nn}{np}

\PY{n}{numero\PYZus{}barra} \PY{o}{=} \PY{n+nb}{list}\PY{p}{(}\PY{n+nb}{range}\PY{p}{(}\PY{l+m+mi}{1}\PY{p}{,} \PY{l+m+mi}{16}\PY{p}{)}\PY{p}{)}

\PY{n}{cargas} \PY{o}{=} \PY{p}{[}
    \PY{l+m+mi}{10350}\PY{p}{,} \PY{l+m+mi}{8450}\PY{p}{,} \PY{l+m+mi}{7200}\PY{p}{,} \PY{l+m+mi}{5100}\PY{p}{,} \PY{l+m+mi}{6500}\PY{p}{,}
    \PY{l+m+mi}{10600}\PY{p}{,} \PY{l+m+mi}{6000}\PY{p}{,} \PY{l+m+mi}{6000}\PY{p}{,} \PY{l+m+mi}{9500}\PY{p}{,} \PY{l+m+mi}{6500}\PY{p}{,}
    \PY{l+m+mi}{9300}\PY{p}{,} \PY{l+m+mi}{6000}\PY{p}{,} \PY{l+m+mi}{6000}\PY{p}{,} \PY{l+m+mi}{5800}\PY{p}{,} \PY{l+m+mi}{6500}
\PY{p}{]}

\PY{n}{tabela\PYZus{}2} \PY{o}{=} \PY{n}{pd}\PY{o}{.}\PY{n}{DataFrame}\PY{p}{(}\PY{p}{\PYZob{}}
    \PY{l+s+s2}{\PYZdq{}}\PY{l+s+s2}{Número da Barra (i)}\PY{l+s+s2}{\PYZdq{}}\PY{p}{:} \PY{n}{numero\PYZus{}barra}\PY{p}{,}
    \PY{l+s+s2}{\PYZdq{}}\PY{l+s+s2}{Carga primeira rachadura (xi)}\PY{l+s+s2}{\PYZdq{}}\PY{p}{:} \PY{n}{cargas}
\PY{p}{\PYZcb{}}\PY{p}{)}

\PY{n}{tabela\PYZus{}2} \PY{o}{=} \PY{p}{(}
    \PY{n}{tabela\PYZus{}2}
    \PY{o}{.}\PY{n}{sort\PYZus{}values}\PY{p}{(}\PY{n}{by}\PY{o}{=}\PY{l+s+s2}{\PYZdq{}}\PY{l+s+s2}{Carga primeira rachadura (xi)}\PY{l+s+s2}{\PYZdq{}}\PY{p}{)}
    \PY{o}{.}\PY{n}{reset\PYZus{}index}\PY{p}{(}\PY{n}{drop}\PY{o}{=}\PY{k+kc}{True}\PY{p}{)}
\PY{p}{)}

\PY{n}{tabela\PYZus{}2}
\end{Verbatim}
\end{tcolorbox}

            \begin{tcolorbox}[breakable, size=fbox, boxrule=.5pt, pad at break*=1mm, opacityfill=0]
\prompt{Out}{outcolor}{3}{\boxspacing}
\begin{Verbatim}[commandchars=\\\{\}]
    Número da Barra (i)  Carga primeira rachadura (xi)
0                     4                           5100
1                    14                           5800
2                     7                           6000
3                     8                           6000
4                    12                           6000
5                    13                           6000
6                    10                           6500
7                     5                           6500
8                    15                           6500
9                     3                           7200
10                    2                           8450
11                   11                           9300
12                    9                           9500
13                    1                          10350
14                    6                          10600
\end{Verbatim}
\end{tcolorbox}
        
    Dividindo em classes de 1000 na carga teremos:

{\def\LTcaptype{table} % do not increment counter
\begin{longtable}[]{@{}rr@{}}
\toprule\noalign{}
Nº da classe & frequência \\
\midrule\noalign{}
\endhead
\bottomrule\noalign{}
\endlastfoot
1 & 5100 a 6100 \\
2 & 6100 a 7100 \\
3 & 7100 a 8100 \\
4 & 8100 a 9100 \\
5 & 9100 a 10100 \\
6 & 10100 a 10600 \\
\end{longtable}
}

ou ainda pode-fazer pela fórmula de Sturges, onde:

\[k=1+3,332\cdot \log(N)\]

onde \(N\) é o número de observações. Então:

o valor mínimo é 5100 e; o valor máximo é 10600.

    \begin{tcolorbox}[breakable, size=fbox, boxrule=1pt, pad at break*=1mm,colback=cellbackground, colframe=cellborder]
\prompt{In}{incolor}{4}{\boxspacing}
\begin{Verbatim}[commandchars=\\\{\}]
\PY{n}{N} \PY{o}{=} \PY{l+m+mi}{15}
\PY{n}{k} \PY{o}{=} \PY{l+m+mi}{1} \PY{o}{+} \PY{l+m+mf}{3.332}\PY{o}{*}\PY{n}{np}\PY{o}{.}\PY{n}{log10}\PY{p}{(}\PY{n}{N}\PY{p}{)}
\PY{n}{k} \PY{o}{=} \PY{l+m+mi}{5}
\PY{n+nb}{print}\PY{p}{(}\PY{l+s+sa}{f}\PY{l+s+s1}{\PYZsq{}}\PY{l+s+s1}{O número de classes é: }\PY{l+s+si}{\PYZob{}}\PY{n}{k}\PY{l+s+si}{\PYZcb{}}\PY{l+s+s1}{\PYZsq{}}\PY{p}{)}
\end{Verbatim}
\end{tcolorbox}

    \begin{Verbatim}[commandchars=\\\{\}]
O número de classes é: 5
    \end{Verbatim}

    \begin{tcolorbox}[breakable, size=fbox, boxrule=1pt, pad at break*=1mm,colback=cellbackground, colframe=cellborder]
\prompt{In}{incolor}{5}{\boxspacing}
\begin{Verbatim}[commandchars=\\\{\}]
\PY{n}{maximo} \PY{o}{=} \PY{l+m+mi}{10600}
\PY{n}{minimo} \PY{o}{=} \PY{l+m+mi}{5100}

\PY{n}{amplitude} \PY{o}{=} \PY{n}{maximo} \PY{o}{\PYZhy{}} \PY{n}{minimo}
\PY{n}{amplitude\PYZus{}classe} \PY{o}{=} \PY{n}{amplitude} \PY{o}{/} \PY{l+m+mi}{5}
\PY{n+nb}{print}\PY{p}{(}\PY{l+s+sa}{f}\PY{l+s+s1}{\PYZsq{}}\PY{l+s+s1}{A Amplitude é igual: }\PY{l+s+si}{\PYZob{}}\PY{n}{amplitude}\PY{l+s+si}{\PYZcb{}}\PY{l+s+s1}{\PYZsq{}}\PY{p}{)}
\PY{n+nb}{print}\PY{p}{(}\PY{l+s+sa}{f}\PY{l+s+s1}{\PYZsq{}}\PY{l+s+s1}{A Amplitude da classe é igual: }\PY{l+s+si}{\PYZob{}}\PY{n}{amplitude\PYZus{}classe}\PY{l+s+si}{\PYZcb{}}\PY{l+s+s1}{\PYZsq{}}\PY{p}{)}
\end{Verbatim}
\end{tcolorbox}

    \begin{Verbatim}[commandchars=\\\{\}]
A Amplitude é igual: 5500
A Amplitude da classe é igual: 1100.0
    \end{Verbatim}

    \begin{tcolorbox}[breakable, size=fbox, boxrule=1pt, pad at break*=1mm,colback=cellbackground, colframe=cellborder]
\prompt{In}{incolor}{6}{\boxspacing}
\begin{Verbatim}[commandchars=\\\{\}]
\PY{n}{classes\PYZus{}carga} \PY{o}{=} \PY{p}{[}
                \PY{l+s+s2}{\PYZdq{}}\PY{l+s+s2}{5100 a 6200}\PY{l+s+s2}{\PYZdq{}}\PY{p}{,}
                \PY{l+s+s2}{\PYZdq{}}\PY{l+s+s2}{6200 a 7300}\PY{l+s+s2}{\PYZdq{}}\PY{p}{,}
                \PY{l+s+s2}{\PYZdq{}}\PY{l+s+s2}{7300 a 8400}\PY{l+s+s2}{\PYZdq{}}\PY{p}{,}
                \PY{l+s+s2}{\PYZdq{}}\PY{l+s+s2}{8400 a 9500}\PY{l+s+s2}{\PYZdq{}}\PY{p}{,}
                \PY{l+s+s2}{\PYZdq{}}\PY{l+s+s2}{9500 a 10600}\PY{l+s+s2}{\PYZdq{}}
                \PY{p}{]}

\PY{n}{carga} \PY{o}{=} \PY{n}{tabela\PYZus{}2}\PY{p}{[}\PY{l+s+s2}{\PYZdq{}}\PY{l+s+s2}{Carga primeira rachadura (xi)}\PY{l+s+s2}{\PYZdq{}}\PY{p}{]}

\PY{n}{soma\PYZus{}f1} \PY{o}{=} \PY{n}{carga}\PY{p}{[}
    \PY{p}{(}\PY{n}{carga} \PY{o}{\PYZgt{}}\PY{o}{=} \PY{l+m+mi}{5100}\PY{p}{)} \PY{o}{\PYZam{}} \PY{p}{(}\PY{n}{carga} \PY{o}{\PYZlt{}} \PY{l+m+mi}{6200}\PY{p}{)}
\PY{p}{]}\PY{o}{.}\PY{n}{count}\PY{p}{(}\PY{p}{)}

\PY{n}{soma\PYZus{}f2} \PY{o}{=} \PY{n}{carga}\PY{p}{[}
    \PY{p}{(}\PY{n}{carga} \PY{o}{\PYZgt{}}\PY{o}{=} \PY{l+m+mi}{6200}\PY{p}{)} \PY{o}{\PYZam{}} \PY{p}{(}\PY{n}{carga} \PY{o}{\PYZlt{}} \PY{l+m+mi}{7300}\PY{p}{)}
\PY{p}{]}\PY{o}{.}\PY{n}{count}\PY{p}{(}\PY{p}{)}

\PY{n}{soma\PYZus{}f3} \PY{o}{=} \PY{n}{carga}\PY{p}{[}
    \PY{p}{(}\PY{n}{carga} \PY{o}{\PYZgt{}}\PY{o}{=} \PY{l+m+mi}{7300}\PY{p}{)} \PY{o}{\PYZam{}} \PY{p}{(}\PY{n}{carga} \PY{o}{\PYZlt{}} \PY{l+m+mi}{8400}\PY{p}{)}
\PY{p}{]}\PY{o}{.}\PY{n}{count}\PY{p}{(}\PY{p}{)}

\PY{n}{soma\PYZus{}f4} \PY{o}{=} \PY{n}{carga}\PY{p}{[}
    \PY{p}{(}\PY{n}{carga} \PY{o}{\PYZgt{}}\PY{o}{=} \PY{l+m+mi}{8400}\PY{p}{)} \PY{o}{\PYZam{}} \PY{p}{(}\PY{n}{carga} \PY{o}{\PYZlt{}} \PY{l+m+mi}{9500}\PY{p}{)}
\PY{p}{]}\PY{o}{.}\PY{n}{count}\PY{p}{(}\PY{p}{)}

\PY{n}{soma\PYZus{}f5} \PY{o}{=} \PY{n}{carga}\PY{p}{[}
    \PY{p}{(}\PY{n}{carga} \PY{o}{\PYZgt{}}\PY{o}{=} \PY{l+m+mi}{9500}\PY{p}{)} \PY{o}{\PYZam{}} \PY{p}{(}\PY{n}{carga} \PY{o}{\PYZlt{}}\PY{o}{=} \PY{l+m+mi}{10600}\PY{p}{)}
\PY{p}{]}\PY{o}{.}\PY{n}{count}\PY{p}{(}\PY{p}{)}


\PY{n}{frequencias\PYZus{}classes} \PY{o}{=} \PY{p}{[}
    \PY{n}{soma\PYZus{}f1}\PY{p}{,}
    \PY{n}{soma\PYZus{}f2}\PY{p}{,}
    \PY{n}{soma\PYZus{}f3}\PY{p}{,}
    \PY{n}{soma\PYZus{}f4}\PY{p}{,}
    \PY{n}{soma\PYZus{}f5}
\PY{p}{]}

\PY{n}{dados\PYZus{}carga} \PY{o}{=} \PY{n}{pd}\PY{o}{.}\PY{n}{DataFrame}\PY{p}{(}\PY{p}{\PYZob{}}
    \PY{l+s+s2}{\PYZdq{}}\PY{l+s+s2}{Classes de Carga}\PY{l+s+s2}{\PYZdq{}}\PY{p}{:} \PY{n}{classes\PYZus{}carga}\PY{p}{,}
    \PY{l+s+s2}{\PYZdq{}}\PY{l+s+s2}{fi}\PY{l+s+s2}{\PYZdq{}}\PY{p}{:} \PY{n}{frequencias\PYZus{}classes}
\PY{p}{\PYZcb{}}\PY{p}{)}

\PY{n}{n} \PY{o}{=} \PY{n}{dados\PYZus{}carga}\PY{p}{[}\PY{l+s+s2}{\PYZdq{}}\PY{l+s+s2}{fi}\PY{l+s+s2}{\PYZdq{}}\PY{p}{]}\PY{o}{.}\PY{n}{sum}\PY{p}{(}\PY{p}{)}

\PY{n}{dados\PYZus{}carga}\PY{p}{[}\PY{l+s+s1}{\PYZsq{}}\PY{l+s+s1}{fri}\PY{l+s+s1}{\PYZsq{}}\PY{p}{]} \PY{o}{=} \PY{n}{dados\PYZus{}carga}\PY{p}{[}\PY{l+s+s1}{\PYZsq{}}\PY{l+s+s1}{fi}\PY{l+s+s1}{\PYZsq{}}\PY{p}{]} \PY{o}{/} \PY{n}{n}

\PY{n}{dados\PYZus{}carga}\PY{p}{[}\PY{l+s+s1}{\PYZsq{}}\PY{l+s+s1}{Fra}\PY{l+s+s1}{\PYZsq{}}\PY{p}{]} \PY{o}{=} \PY{n}{dados\PYZus{}carga}\PY{p}{[}\PY{l+s+s1}{\PYZsq{}}\PY{l+s+s1}{fri}\PY{l+s+s1}{\PYZsq{}}\PY{p}{]}\PY{o}{.}\PY{n}{cumsum}\PY{p}{(}\PY{p}{)}

\PY{n}{centro\PYZus{}classe} \PY{o}{=} \PY{p}{\PYZob{}}\PY{l+s+s1}{\PYZsq{}}\PY{l+s+s1}{Centro da Classe}\PY{l+s+s1}{\PYZsq{}}\PY{p}{:} \PY{p}{[}\PY{p}{(}\PY{l+m+mi}{5100}\PY{o}{+}\PY{l+m+mi}{6200}\PY{p}{)}\PY{o}{/}\PY{l+m+mi}{2}\PY{p}{,} \PY{p}{(}\PY{l+m+mi}{6200}\PY{o}{+}\PY{l+m+mi}{7300}\PY{p}{)}\PY{o}{/}\PY{l+m+mi}{2}\PY{p}{,} \PY{p}{(}\PY{l+m+mi}{7300}\PY{o}{+}\PY{l+m+mi}{8400}\PY{p}{)}\PY{o}{/}\PY{l+m+mi}{2}\PY{p}{,} \PY{p}{(}\PY{l+m+mi}{8400}\PY{o}{+}\PY{l+m+mi}{9500}\PY{p}{)}\PY{o}{/}\PY{l+m+mi}{2}\PY{p}{,} \PY{p}{(}\PY{l+m+mi}{9500}\PY{o}{+}\PY{l+m+mi}{10600}\PY{p}{)}\PY{o}{/}\PY{l+m+mi}{2}\PY{p}{]}\PY{p}{\PYZcb{}}

\PY{n}{dados\PYZus{}carga}\PY{p}{[}\PY{l+s+s1}{\PYZsq{}}\PY{l+s+s1}{Centro da Classe}\PY{l+s+s1}{\PYZsq{}}\PY{p}{]} \PY{o}{=} \PY{n}{centro\PYZus{}classe}\PY{p}{[}\PY{l+s+s1}{\PYZsq{}}\PY{l+s+s1}{Centro da Classe}\PY{l+s+s1}{\PYZsq{}}\PY{p}{]}


\PY{n}{media} \PY{o}{=} \PY{p}{(}
    \PY{n}{dados\PYZus{}carga}\PY{p}{[}\PY{l+s+s1}{\PYZsq{}}\PY{l+s+s1}{fri}\PY{l+s+s1}{\PYZsq{}}\PY{p}{]} \PY{o}{*}
    \PY{n}{dados\PYZus{}carga}\PY{p}{[}\PY{l+s+s1}{\PYZsq{}}\PY{l+s+s1}{Centro da Classe}\PY{l+s+s1}{\PYZsq{}}\PY{p}{]}
\PY{p}{)}\PY{o}{.}\PY{n}{sum}\PY{p}{(}\PY{p}{)}

\PY{n}{var} \PY{o}{=} \PY{p}{(}
    \PY{n}{dados\PYZus{}carga}\PY{p}{[}\PY{l+s+s1}{\PYZsq{}}\PY{l+s+s1}{fri}\PY{l+s+s1}{\PYZsq{}}\PY{p}{]} \PY{o}{*}
    \PY{p}{(}\PY{n}{dados\PYZus{}carga}\PY{p}{[}\PY{l+s+s1}{\PYZsq{}}\PY{l+s+s1}{Centro da Classe}\PY{l+s+s1}{\PYZsq{}}\PY{p}{]} \PY{o}{\PYZhy{}} \PY{n}{media}\PY{p}{)}\PY{o}{*}\PY{o}{*}\PY{l+m+mi}{2}
\PY{p}{)}\PY{o}{.}\PY{n}{sum}\PY{p}{(}\PY{p}{)}

\PY{n}{std} \PY{o}{=} \PY{n}{np}\PY{o}{.}\PY{n}{sqrt}\PY{p}{(}\PY{n}{var}\PY{p}{)}

\PY{n+nb}{print}\PY{p}{(}\PY{l+s+sa}{f}\PY{l+s+s1}{\PYZsq{}}\PY{l+s+s1}{Média agrupada: }\PY{l+s+si}{\PYZob{}}\PY{n}{media}\PY{l+s+si}{:}\PY{l+s+s1}{.2f}\PY{l+s+si}{\PYZcb{}}\PY{l+s+s1}{\PYZsq{}}\PY{p}{)}
\PY{n+nb}{print}\PY{p}{(}\PY{l+s+sa}{f}\PY{l+s+s1}{\PYZsq{}}\PY{l+s+s1}{Variância agrupada: }\PY{l+s+si}{\PYZob{}}\PY{n}{var}\PY{l+s+si}{:}\PY{l+s+s1}{.2f}\PY{l+s+si}{\PYZcb{}}\PY{l+s+s1}{\PYZsq{}}\PY{p}{)}
\PY{n+nb}{print}\PY{p}{(}\PY{l+s+sa}{f}\PY{l+s+s1}{\PYZsq{}}\PY{l+s+s1}{Desvio padrão agrupado: }\PY{l+s+si}{\PYZob{}}\PY{n}{std}\PY{l+s+si}{:}\PY{l+s+s1}{.2f}\PY{l+s+si}{\PYZcb{}}\PY{l+s+s1}{\PYZsq{}}\PY{p}{)}
\end{Verbatim}
\end{tcolorbox}

    \begin{Verbatim}[commandchars=\\\{\}]
Média agrupada: 7263.33
Variância agrupada: 3043822.22
Desvio padrão agrupado: 1744.66
    \end{Verbatim}

    \begin{tcolorbox}[breakable, size=fbox, boxrule=1pt, pad at break*=1mm,colback=cellbackground, colframe=cellborder]
\prompt{In}{incolor}{7}{\boxspacing}
\begin{Verbatim}[commandchars=\\\{\}]
\PY{n}{dados\PYZus{}carga}
\end{Verbatim}
\end{tcolorbox}

            \begin{tcolorbox}[breakable, size=fbox, boxrule=.5pt, pad at break*=1mm, opacityfill=0]
\prompt{Out}{outcolor}{7}{\boxspacing}
\begin{Verbatim}[commandchars=\\\{\}]
  Classes de Carga  fi       fri       Fra  Centro da Classe
0      5100 a 6200   6  0.400000  0.400000            5650.0
1      6200 a 7300   4  0.266667  0.666667            6750.0
2      7300 a 8400   0  0.000000  0.666667            7850.0
3      8400 a 9500   2  0.133333  0.800000            8950.0
4     9500 a 10600   3  0.200000  1.000000           10050.0
\end{Verbatim}
\end{tcolorbox}
        
    \subsection{As cargas de colapso
são:}\label{as-cargas-de-colapso-suxe3o}

{\def\LTcaptype{table} % do not increment counter
\begin{longtable}[]{@{}rr@{}}
\toprule\noalign{}
Nº da barra & Carga de colapso \\
\midrule\noalign{}
\endhead
\bottomrule\noalign{}
\endlastfoot
1 & 10350 \\
2 & 9300 \\
3 & 9600 \\
4 & 10300 \\
5 & 9400 \\
6 & 10600 \\
7 & 10100 \\
8 & 9900 \\
9 & 9500 \\
10 & 10200 \\
11 & 9300 \\
12 & 9550 \\
13 & 9550 \\
14 & 10500 \\
15 & 10200 \\
\end{longtable}
}

    \begin{tcolorbox}[breakable, size=fbox, boxrule=1pt, pad at break*=1mm,colback=cellbackground, colframe=cellborder]
\prompt{In}{incolor}{8}{\boxspacing}
\begin{Verbatim}[commandchars=\\\{\}]
\PY{n}{cargas\PYZus{}colapso} \PY{o}{=} \PY{p}{[}\PY{l+m+mi}{10350}\PY{p}{,} \PY{l+m+mi}{9300}\PY{p}{,} \PY{l+m+mi}{9600}\PY{p}{,} \PY{l+m+mi}{10300}\PY{p}{,} \PY{l+m+mi}{9400}\PY{p}{,}
                  \PY{l+m+mi}{10600}\PY{p}{,} \PY{l+m+mi}{10100}\PY{p}{,} \PY{l+m+mi}{9900}\PY{p}{,} \PY{l+m+mi}{9500}\PY{p}{,} \PY{l+m+mi}{10200}\PY{p}{,}
                  \PY{l+m+mi}{9300}\PY{p}{,} \PY{l+m+mi}{9550}\PY{p}{,} \PY{l+m+mi}{9550}\PY{p}{,} \PY{l+m+mi}{10500}\PY{p}{,} \PY{l+m+mi}{10200}\PY{p}{]}

\PY{n}{tabela\PYZus{}3} \PY{o}{=} \PY{n}{pd}\PY{o}{.}\PY{n}{DataFrame}\PY{p}{(}\PY{p}{\PYZob{}}
    \PY{l+s+s2}{\PYZdq{}}\PY{l+s+s2}{Número da Barra (i)}\PY{l+s+s2}{\PYZdq{}}\PY{p}{:} \PY{n}{numero\PYZus{}barra}\PY{p}{,}
    \PY{l+s+s2}{\PYZdq{}}\PY{l+s+s2}{Carga de Colapso (xi)}\PY{l+s+s2}{\PYZdq{}}\PY{p}{:} \PY{n}{cargas\PYZus{}colapso}
\PY{p}{\PYZcb{}}\PY{p}{)}

\PY{n}{tabela\PYZus{}3} \PY{o}{=} \PY{p}{(}
    \PY{n}{tabela\PYZus{}3}
    \PY{o}{.}\PY{n}{sort\PYZus{}values}\PY{p}{(}\PY{n}{by}\PY{o}{=}\PY{l+s+s2}{\PYZdq{}}\PY{l+s+s2}{Carga de Colapso (xi)}\PY{l+s+s2}{\PYZdq{}}\PY{p}{)}
    \PY{o}{.}\PY{n}{reset\PYZus{}index}\PY{p}{(}\PY{n}{drop}\PY{o}{=}\PY{k+kc}{True}\PY{p}{)}
\PY{p}{)}

\PY{n}{tabela\PYZus{}3}
\end{Verbatim}
\end{tcolorbox}

            \begin{tcolorbox}[breakable, size=fbox, boxrule=.5pt, pad at break*=1mm, opacityfill=0]
\prompt{Out}{outcolor}{8}{\boxspacing}
\begin{Verbatim}[commandchars=\\\{\}]
    Número da Barra (i)  Carga de Colapso (xi)
0                     2                   9300
1                    11                   9300
2                     5                   9400
3                     9                   9500
4                    13                   9550
5                    12                   9550
6                     3                   9600
7                     8                   9900
8                     7                  10100
9                    10                  10200
10                   15                  10200
11                    4                  10300
12                    1                  10350
13                   14                  10500
14                    6                  10600
\end{Verbatim}
\end{tcolorbox}
        
    Sendo o número de observações o mesmo do item anterior, usaremos 5
classes.\\
então as classes seriam dadas por:

    \begin{tcolorbox}[breakable, size=fbox, boxrule=1pt, pad at break*=1mm,colback=cellbackground, colframe=cellborder]
\prompt{In}{incolor}{9}{\boxspacing}
\begin{Verbatim}[commandchars=\\\{\}]
\PY{n}{max\PYZus{}colapso} \PY{o}{=} \PY{n}{tabela\PYZus{}3}\PY{p}{[}\PY{l+s+s1}{\PYZsq{}}\PY{l+s+s1}{Carga de Colapso (xi)}\PY{l+s+s1}{\PYZsq{}}\PY{p}{]}\PY{o}{.}\PY{n}{max}\PY{p}{(}\PY{p}{)}
\PY{n}{min\PYZus{}colapso} \PY{o}{=} \PY{n}{tabela\PYZus{}3}\PY{p}{[}\PY{l+s+s1}{\PYZsq{}}\PY{l+s+s1}{Carga de Colapso (xi)}\PY{l+s+s1}{\PYZsq{}}\PY{p}{]}\PY{o}{.}\PY{n}{min}\PY{p}{(}\PY{p}{)}
\PY{n}{amplitude\PYZus{}colapso} \PY{o}{=} \PY{n}{max\PYZus{}colapso} \PY{o}{\PYZhy{}} \PY{n}{min\PYZus{}colapso}
\PY{n+nb}{print}\PY{p}{(}\PY{l+s+sa}{f}\PY{l+s+s1}{\PYZsq{}}\PY{l+s+s1}{A amplitude\PYZus{}colapso é: }\PY{l+s+si}{\PYZob{}}\PY{n}{amplitude\PYZus{}colapso}\PY{l+s+si}{:}\PY{l+s+s1}{.2f}\PY{l+s+si}{\PYZcb{}}\PY{l+s+s1}{\PYZsq{}}\PY{p}{)}
\end{Verbatim}
\end{tcolorbox}

    \begin{Verbatim}[commandchars=\\\{\}]
A amplitude\_colapso é: 1300.00
    \end{Verbatim}

    \begin{tcolorbox}[breakable, size=fbox, boxrule=1pt, pad at break*=1mm,colback=cellbackground, colframe=cellborder]
\prompt{In}{incolor}{10}{\boxspacing}
\begin{Verbatim}[commandchars=\\\{\}]
\PY{n}{classes\PYZus{}carga\PYZus{}colapso} \PY{o}{=} \PY{p}{[}
    \PY{l+s+s2}{\PYZdq{}}\PY{l+s+s2}{9300 a 9560}\PY{l+s+s2}{\PYZdq{}}\PY{p}{,}
    \PY{l+s+s2}{\PYZdq{}}\PY{l+s+s2}{9560 a 9820}\PY{l+s+s2}{\PYZdq{}}\PY{p}{,}
    \PY{l+s+s2}{\PYZdq{}}\PY{l+s+s2}{9820 a 10080}\PY{l+s+s2}{\PYZdq{}}\PY{p}{,}
    \PY{l+s+s2}{\PYZdq{}}\PY{l+s+s2}{10080 a 10340}\PY{l+s+s2}{\PYZdq{}}\PY{p}{,}
    \PY{l+s+s2}{\PYZdq{}}\PY{l+s+s2}{10340 a 10600}\PY{l+s+s2}{\PYZdq{}}
\PY{p}{]}

\PY{n}{carga\PYZus{}colapso} \PY{o}{=} \PY{n}{tabela\PYZus{}3}\PY{p}{[}\PY{l+s+s2}{\PYZdq{}}\PY{l+s+s2}{Carga de Colapso (xi)}\PY{l+s+s2}{\PYZdq{}}\PY{p}{]}

\PY{n}{soma\PYZus{}f1} \PY{o}{=} \PY{n}{carga\PYZus{}colapso}\PY{p}{[}
    \PY{p}{(}\PY{n}{carga\PYZus{}colapso} \PY{o}{\PYZgt{}}\PY{o}{=} \PY{l+m+mi}{9300}\PY{p}{)} \PY{o}{\PYZam{}} \PY{p}{(}\PY{n}{carga\PYZus{}colapso} \PY{o}{\PYZlt{}} \PY{l+m+mi}{9560}\PY{p}{)}
\PY{p}{]}\PY{o}{.}\PY{n}{count}\PY{p}{(}\PY{p}{)}

\PY{n}{soma\PYZus{}f2} \PY{o}{=} \PY{n}{carga\PYZus{}colapso}\PY{p}{[}
    \PY{p}{(}\PY{n}{carga\PYZus{}colapso} \PY{o}{\PYZgt{}}\PY{o}{=} \PY{l+m+mi}{9560}\PY{p}{)} \PY{o}{\PYZam{}} \PY{p}{(}\PY{n}{carga\PYZus{}colapso} \PY{o}{\PYZlt{}} \PY{l+m+mi}{9820}\PY{p}{)}
\PY{p}{]}\PY{o}{.}\PY{n}{count}\PY{p}{(}\PY{p}{)}

\PY{n}{soma\PYZus{}f3} \PY{o}{=} \PY{n}{carga\PYZus{}colapso}\PY{p}{[}
    \PY{p}{(}\PY{n}{carga\PYZus{}colapso} \PY{o}{\PYZgt{}}\PY{o}{=} \PY{l+m+mi}{9820}\PY{p}{)} \PY{o}{\PYZam{}} \PY{p}{(}\PY{n}{carga\PYZus{}colapso} \PY{o}{\PYZlt{}} \PY{l+m+mi}{10080}\PY{p}{)}
\PY{p}{]}\PY{o}{.}\PY{n}{count}\PY{p}{(}\PY{p}{)}

\PY{n}{soma\PYZus{}f4} \PY{o}{=} \PY{n}{carga\PYZus{}colapso}\PY{p}{[}
    \PY{p}{(}\PY{n}{carga\PYZus{}colapso} \PY{o}{\PYZgt{}}\PY{o}{=} \PY{l+m+mi}{10080}\PY{p}{)} \PY{o}{\PYZam{}} \PY{p}{(}\PY{n}{carga\PYZus{}colapso} \PY{o}{\PYZlt{}} \PY{l+m+mi}{10340}\PY{p}{)}
\PY{p}{]}\PY{o}{.}\PY{n}{count}\PY{p}{(}\PY{p}{)}

\PY{n}{soma\PYZus{}f5} \PY{o}{=} \PY{n}{carga\PYZus{}colapso}\PY{p}{[}
    \PY{p}{(}\PY{n}{carga\PYZus{}colapso} \PY{o}{\PYZgt{}}\PY{o}{=} \PY{l+m+mi}{10340}\PY{p}{)} \PY{o}{\PYZam{}} \PY{p}{(}\PY{n}{carga\PYZus{}colapso} \PY{o}{\PYZlt{}}\PY{o}{=} \PY{l+m+mi}{10600}\PY{p}{)}
\PY{p}{]}\PY{o}{.}\PY{n}{count}\PY{p}{(}\PY{p}{)}

\PY{n}{frequencias\PYZus{}classes} \PY{o}{=} \PY{p}{[}
    \PY{n}{soma\PYZus{}f1}\PY{p}{,}
    \PY{n}{soma\PYZus{}f2}\PY{p}{,}
    \PY{n}{soma\PYZus{}f3}\PY{p}{,}
    \PY{n}{soma\PYZus{}f4}\PY{p}{,}
    \PY{n}{soma\PYZus{}f5}
\PY{p}{]}

\PY{n+nb}{print}\PY{p}{(}\PY{l+s+s2}{\PYZdq{}}\PY{l+s+s2}{Frequências:}\PY{l+s+s2}{\PYZdq{}}\PY{p}{,} \PY{n}{frequencias\PYZus{}classes}\PY{p}{)}
\PY{n+nb}{print}\PY{p}{(}\PY{l+s+s2}{\PYZdq{}}\PY{l+s+s2}{Total contado:}\PY{l+s+s2}{\PYZdq{}}\PY{p}{,} \PY{n+nb}{sum}\PY{p}{(}\PY{n}{frequencias\PYZus{}classes}\PY{p}{)}\PY{p}{)}
\PY{n+nb}{print}\PY{p}{(}\PY{l+s+s2}{\PYZdq{}}\PY{l+s+s2}{Total original:}\PY{l+s+s2}{\PYZdq{}}\PY{p}{,} \PY{n+nb}{len}\PY{p}{(}\PY{n}{carga\PYZus{}colapso}\PY{p}{)}\PY{p}{)}

\PY{n}{dados\PYZus{}carga\PYZus{}colapso} \PY{o}{=} \PY{n}{pd}\PY{o}{.}\PY{n}{DataFrame}\PY{p}{(}\PY{p}{\PYZob{}}
    \PY{l+s+s2}{\PYZdq{}}\PY{l+s+s2}{Classes de Carga}\PY{l+s+s2}{\PYZdq{}}\PY{p}{:} \PY{n}{classes\PYZus{}carga\PYZus{}colapso}\PY{p}{,}
    \PY{l+s+s2}{\PYZdq{}}\PY{l+s+s2}{fi}\PY{l+s+s2}{\PYZdq{}}\PY{p}{:} \PY{n}{frequencias\PYZus{}classes}
\PY{p}{\PYZcb{}}\PY{p}{)}

\PY{n}{n} \PY{o}{=} \PY{n}{dados\PYZus{}carga\PYZus{}colapso}\PY{p}{[}\PY{l+s+s2}{\PYZdq{}}\PY{l+s+s2}{fi}\PY{l+s+s2}{\PYZdq{}}\PY{p}{]}\PY{o}{.}\PY{n}{sum}\PY{p}{(}\PY{p}{)}

\PY{n}{centro\PYZus{}classe\PYZus{}colapso} \PY{o}{=} \PY{p}{[}
    \PY{p}{(}\PY{l+m+mi}{9300} \PY{o}{+} \PY{l+m+mi}{9560}\PY{p}{)} \PY{o}{/} \PY{l+m+mi}{2}\PY{p}{,}
    \PY{p}{(}\PY{l+m+mi}{9560} \PY{o}{+} \PY{l+m+mi}{9820}\PY{p}{)} \PY{o}{/} \PY{l+m+mi}{2}\PY{p}{,}
    \PY{p}{(}\PY{l+m+mi}{9820} \PY{o}{+} \PY{l+m+mi}{10080}\PY{p}{)} \PY{o}{/} \PY{l+m+mi}{2}\PY{p}{,}
    \PY{p}{(}\PY{l+m+mi}{10080} \PY{o}{+} \PY{l+m+mi}{10340}\PY{p}{)} \PY{o}{/} \PY{l+m+mi}{2}\PY{p}{,}
    \PY{p}{(}\PY{l+m+mi}{10340} \PY{o}{+} \PY{l+m+mi}{10600}\PY{p}{)} \PY{o}{/} \PY{l+m+mi}{2}
\PY{p}{]}

\PY{n}{dados\PYZus{}carga\PYZus{}colapso}\PY{p}{[}\PY{l+s+s2}{\PYZdq{}}\PY{l+s+s2}{fri}\PY{l+s+s2}{\PYZdq{}}\PY{p}{]} \PY{o}{=} \PY{n}{dados\PYZus{}carga\PYZus{}colapso}\PY{p}{[}\PY{l+s+s2}{\PYZdq{}}\PY{l+s+s2}{fi}\PY{l+s+s2}{\PYZdq{}}\PY{p}{]} \PY{o}{/} \PY{n}{n}
\PY{n}{dados\PYZus{}carga\PYZus{}colapso}\PY{p}{[}\PY{l+s+s2}{\PYZdq{}}\PY{l+s+s2}{Fra}\PY{l+s+s2}{\PYZdq{}}\PY{p}{]} \PY{o}{=} \PY{n}{dados\PYZus{}carga\PYZus{}colapso}\PY{p}{[}\PY{l+s+s2}{\PYZdq{}}\PY{l+s+s2}{fri}\PY{l+s+s2}{\PYZdq{}}\PY{p}{]}\PY{o}{.}\PY{n}{cumsum}\PY{p}{(}\PY{p}{)}
\PY{n}{dados\PYZus{}carga\PYZus{}colapso}\PY{p}{[}\PY{l+s+s2}{\PYZdq{}}\PY{l+s+s2}{Centro da Classe}\PY{l+s+s2}{\PYZdq{}}\PY{p}{]} \PY{o}{=} \PY{n}{centro\PYZus{}classe\PYZus{}colapso}

\PY{n}{dados\PYZus{}carga\PYZus{}colapso}
\end{Verbatim}
\end{tcolorbox}

    \begin{Verbatim}[commandchars=\\\{\}]
Frequências: [np.int64(6), np.int64(1), np.int64(1), np.int64(4), np.int64(3)]
Total contado: 15
Total original: 15
    \end{Verbatim}

            \begin{tcolorbox}[breakable, size=fbox, boxrule=.5pt, pad at break*=1mm, opacityfill=0]
\prompt{Out}{outcolor}{10}{\boxspacing}
\begin{Verbatim}[commandchars=\\\{\}]
  Classes de Carga  fi       fri       Fra  Centro da Classe
0      9300 a 9560   6  0.400000  0.400000            9430.0
1      9560 a 9820   1  0.066667  0.466667            9690.0
2     9820 a 10080   1  0.066667  0.533333            9950.0
3    10080 a 10340   4  0.266667  0.800000           10210.0
4    10340 a 10600   3  0.200000  1.000000           10470.0
\end{Verbatim}
\end{tcolorbox}
        
    \subsubsection{Calcular a média, VAR e
std.}\label{calcular-a-muxe9dia-var-e-std.}

    \begin{tcolorbox}[breakable, size=fbox, boxrule=1pt, pad at break*=1mm,colback=cellbackground, colframe=cellborder]
\prompt{In}{incolor}{11}{\boxspacing}
\begin{Verbatim}[commandchars=\\\{\}]
\PY{n}{media\PYZus{}colapso} \PY{o}{=} \PY{p}{(}\PY{n}{dados\PYZus{}carga\PYZus{}colapso}\PY{p}{[}\PY{l+s+s1}{\PYZsq{}}\PY{l+s+s1}{fri}\PY{l+s+s1}{\PYZsq{}}\PY{p}{]} \PY{o}{*} \PY{n}{dados\PYZus{}carga\PYZus{}colapso}\PY{p}{[}\PY{l+s+s1}{\PYZsq{}}\PY{l+s+s1}{Centro da Classe}\PY{l+s+s1}{\PYZsq{}}\PY{p}{]}\PY{p}{)}\PY{o}{.}\PY{n}{sum}\PY{p}{(}\PY{p}{)}
\PY{n}{var\PYZus{}colapso} \PY{o}{=} \PY{p}{(}\PY{p}{(}\PY{p}{(}\PY{n}{dados\PYZus{}carga\PYZus{}colapso}\PY{p}{[}\PY{l+s+s1}{\PYZsq{}}\PY{l+s+s1}{Centro da Classe}\PY{l+s+s1}{\PYZsq{}}\PY{p}{]}\PY{o}{\PYZhy{}}\PY{n}{media\PYZus{}colapso}\PY{p}{)}\PY{o}{*}\PY{o}{*}\PY{l+m+mi}{2}\PY{p}{)} \PY{o}{*} \PY{n}{dados\PYZus{}carga\PYZus{}colapso}\PY{p}{[}\PY{l+s+s1}{\PYZsq{}}\PY{l+s+s1}{fri}\PY{l+s+s1}{\PYZsq{}}\PY{p}{]}\PY{p}{)}\PY{o}{.}\PY{n}{sum}\PY{p}{(}\PY{p}{)}
\PY{n}{std\PYZus{}colapso} \PY{o}{=} \PY{n}{np}\PY{o}{.}\PY{n}{sqrt}\PY{p}{(}\PY{n}{var\PYZus{}colapso}\PY{p}{)}
\PY{n+nb}{print}\PY{p}{(}\PY{n}{media\PYZus{}colapso}\PY{p}{)}
\PY{n+nb}{print}\PY{p}{(}\PY{n}{var\PYZus{}colapso}\PY{p}{)}
\PY{n+nb}{print}\PY{p}{(}\PY{n}{std\PYZus{}colapso}\PY{p}{)}
\end{Verbatim}
\end{tcolorbox}

    \begin{Verbatim}[commandchars=\\\{\}]
9898.0
182069.3333333333
426.6958323364939
    \end{Verbatim}

    3-- Sejam os dados de tempos de reação a um estímulo (Y) e acuidade
visual (Z) de 20 indivíduos.

{\def\LTcaptype{table} % do not increment counter
\begin{longtable}[]{@{}rrr@{}}
\toprule\noalign{}
Indivíduo & Y & Z \\
\midrule\noalign{}
\endhead
\bottomrule\noalign{}
\endlastfoot
1 & 96 & 90 \\
2 & 92 & 100 \\
3 & 106 & 80 \\
4 & 100 & 90 \\
5 & 98 & 100 \\
6 & 104 & 90 \\
7 & 110 & 80 \\
8 & 101 & 90 \\
9 & 116 & 70 \\
10 & 106 & 90 \\
11 & 109 & 90 \\
12 & 100 & 80 \\
13 & 112 & 90 \\
14 & 105 & 80 \\
15 & 118 & 70 \\
16 & 108 & 90 \\
17 & 113 & 90 \\
18 & 112 & 90 \\
19 & 127 & 60 \\
20 & 117 & 80 \\
\end{longtable}
}

\begin{itemize}
\tightlist
\item
  Calcule a distribuição de frequências das duas variáveis aleatórias
  (considere a representação individual ou por classe).
\item
  Calcule a distribuição de frequências conjunta das duas variáveis
  aleatórias (considere a representação individual).\\
\item
  Calcule o coeficiente de variação de cada uma das duas variáveis
  aleatórias e interprete-os.\\
\item
  Qual a frequência de tempo de reação a um estímulo maior que 100? Qual
  a frequência de tempo de reação a um estímulo menor que 120? Qual a
  frequência de tempo de reação a um estímulo maior que 100 e menor que
  120?\\
\item
  Calcule os quartis das duas distribuições de frequências. Monte o
  diagrama de Box-Plot.\\
\item
  Calcule o coeficiente de correlação entre os tempos de reação a um
  estímulo e a acuidade visual desse conjunto de indivíduos.
  Interprete-o.
\end{itemize}

    \subsubsection{Para a variável y}\label{para-a-variuxe1vel-y}

    \begin{tcolorbox}[breakable, size=fbox, boxrule=1pt, pad at break*=1mm,colback=cellbackground, colframe=cellborder]
\prompt{In}{incolor}{12}{\boxspacing}
\begin{Verbatim}[commandchars=\\\{\}]
\PY{c+c1}{\PYZsh{} distribuição de frequências}
\PY{k+kn}{import}\PY{+w}{ }\PY{n+nn}{matplotlib}\PY{n+nn}{.}\PY{n+nn}{pyplot}\PY{+w}{ }\PY{k}{as}\PY{+w}{ }\PY{n+nn}{plt}
\PY{n}{n} \PY{o}{=} \PY{l+m+mi}{20}
\PY{n}{individuo} \PY{o}{=} \PY{n+nb}{list}\PY{p}{(}\PY{n+nb}{range}\PY{p}{(}\PY{l+m+mi}{1}\PY{p}{,} \PY{l+m+mi}{21}\PY{p}{)}\PY{p}{)}

\PY{n}{y} \PY{o}{=} \PY{p}{[}\PY{l+m+mi}{96}\PY{p}{,} \PY{l+m+mi}{92}\PY{p}{,} \PY{l+m+mi}{106}\PY{p}{,} \PY{l+m+mi}{100}\PY{p}{,} \PY{l+m+mi}{98}\PY{p}{,} \PY{l+m+mi}{104}\PY{p}{,} \PY{l+m+mi}{110}\PY{p}{,} \PY{l+m+mi}{101}\PY{p}{,} \PY{l+m+mi}{116}\PY{p}{,} \PY{l+m+mi}{106}\PY{p}{,} \PY{l+m+mi}{109}\PY{p}{,} \PY{l+m+mi}{100}\PY{p}{,} \PY{l+m+mi}{112}\PY{p}{,} \PY{l+m+mi}{105}\PY{p}{,} \PY{l+m+mi}{118}\PY{p}{,} \PY{l+m+mi}{108}\PY{p}{,} \PY{l+m+mi}{113}\PY{p}{,} \PY{l+m+mi}{112}\PY{p}{,} \PY{l+m+mi}{127}\PY{p}{,} \PY{l+m+mi}{117}\PY{p}{]}
\PY{n}{z} \PY{o}{=} \PY{p}{[}\PY{l+m+mi}{90}\PY{p}{,} \PY{l+m+mi}{100}\PY{p}{,} \PY{l+m+mi}{80}\PY{p}{,} \PY{l+m+mi}{90}\PY{p}{,} \PY{l+m+mi}{100}\PY{p}{,} \PY{l+m+mi}{90}\PY{p}{,} \PY{l+m+mi}{80}\PY{p}{,} \PY{l+m+mi}{90}\PY{p}{,} \PY{l+m+mi}{70}\PY{p}{,} \PY{l+m+mi}{100}\PY{p}{,} \PY{l+m+mi}{90}\PY{p}{,} \PY{l+m+mi}{80}\PY{p}{,} \PY{l+m+mi}{90}\PY{p}{,} \PY{l+m+mi}{80}\PY{p}{,} \PY{l+m+mi}{70}\PY{p}{,} \PY{l+m+mi}{90}\PY{p}{,} \PY{l+m+mi}{90}\PY{p}{,} \PY{l+m+mi}{90}\PY{p}{,} \PY{l+m+mi}{60}\PY{p}{,} \PY{l+m+mi}{80}\PY{p}{]}
\PY{n}{df} \PY{o}{=} \PY{n}{pd}\PY{o}{.}\PY{n}{DataFrame}\PY{p}{(}\PY{p}{\PYZob{}}
    \PY{l+s+s2}{\PYZdq{}}\PY{l+s+s2}{indivíduo}\PY{l+s+s2}{\PYZdq{}}\PY{p}{:} \PY{n}{individuo}\PY{p}{,}
    \PY{l+s+s2}{\PYZdq{}}\PY{l+s+s2}{y}\PY{l+s+s2}{\PYZdq{}}\PY{p}{:} \PY{n}{y}\PY{p}{,}
    \PY{l+s+s2}{\PYZdq{}}\PY{l+s+s2}{z}\PY{l+s+s2}{\PYZdq{}}\PY{p}{:} \PY{n}{z}
\PY{p}{\PYZcb{}}\PY{p}{)}
\PY{n}{freq\PYZus{}y} \PY{o}{=} \PY{p}{(}
    \PY{n}{df}\PY{p}{[}\PY{l+s+s2}{\PYZdq{}}\PY{l+s+s2}{y}\PY{l+s+s2}{\PYZdq{}}\PY{p}{]}
    \PY{o}{.}\PY{n}{value\PYZus{}counts}\PY{p}{(}\PY{p}{)}
    \PY{o}{.}\PY{n}{sort\PYZus{}index}\PY{p}{(}\PY{p}{)}
    \PY{o}{.}\PY{n}{reset\PYZus{}index}\PY{p}{(}\PY{p}{)}
\PY{p}{)}

\PY{n}{freq\PYZus{}y}\PY{o}{.}\PY{n}{columns} \PY{o}{=} \PY{p}{[}\PY{l+s+s2}{\PYZdq{}}\PY{l+s+s2}{Y}\PY{l+s+s2}{\PYZdq{}}\PY{p}{,} \PY{l+s+s2}{\PYZdq{}}\PY{l+s+s2}{fi\PYZus{}y}\PY{l+s+s2}{\PYZdq{}}\PY{p}{]}
\PY{n}{freq\PYZus{}y}\PY{p}{[}\PY{l+s+s1}{\PYZsq{}}\PY{l+s+s1}{Fi\PYZus{}y}\PY{l+s+s1}{\PYZsq{}}\PY{p}{]} \PY{o}{=} \PY{n}{freq\PYZus{}y}\PY{p}{[}\PY{l+s+s1}{\PYZsq{}}\PY{l+s+s1}{fi\PYZus{}y}\PY{l+s+s1}{\PYZsq{}}\PY{p}{]}\PY{o}{.}\PY{n}{cumsum}\PY{p}{(}\PY{p}{)}
\PY{n}{freq\PYZus{}y}\PY{p}{[}\PY{l+s+s1}{\PYZsq{}}\PY{l+s+s1}{fr\PYZus{}y}\PY{l+s+s1}{\PYZsq{}}\PY{p}{]} \PY{o}{=} \PY{n}{freq\PYZus{}y}\PY{p}{[}\PY{l+s+s2}{\PYZdq{}}\PY{l+s+s2}{fi\PYZus{}y}\PY{l+s+s2}{\PYZdq{}}\PY{p}{]} \PY{o}{/} \PY{n}{n}
\PY{n}{freq\PYZus{}y}\PY{p}{[}\PY{l+s+s1}{\PYZsq{}}\PY{l+s+s1}{Frac}\PY{l+s+s1}{\PYZsq{}}\PY{p}{]} \PY{o}{=} \PY{n}{freq\PYZus{}y}\PY{p}{[}\PY{l+s+s2}{\PYZdq{}}\PY{l+s+s2}{fr\PYZus{}y}\PY{l+s+s2}{\PYZdq{}}\PY{p}{]}\PY{o}{.}\PY{n}{cumsum}\PY{p}{(}\PY{p}{)}
\PY{n}{freq\PYZus{}y}
\end{Verbatim}
\end{tcolorbox}

            \begin{tcolorbox}[breakable, size=fbox, boxrule=.5pt, pad at break*=1mm, opacityfill=0]
\prompt{Out}{outcolor}{12}{\boxspacing}
\begin{Verbatim}[commandchars=\\\{\}]
      Y  fi\_y  Fi\_y  fr\_y  Frac
0    92     1     1  0.05  0.05
1    96     1     2  0.05  0.10
2    98     1     3  0.05  0.15
3   100     2     5  0.10  0.25
4   101     1     6  0.05  0.30
5   104     1     7  0.05  0.35
6   105     1     8  0.05  0.40
7   106     2    10  0.10  0.50
8   108     1    11  0.05  0.55
9   109     1    12  0.05  0.60
10  110     1    13  0.05  0.65
11  112     2    15  0.10  0.75
12  113     1    16  0.05  0.80
13  116     1    17  0.05  0.85
14  117     1    18  0.05  0.90
15  118     1    19  0.05  0.95
16  127     1    20  0.05  1.00
\end{Verbatim}
\end{tcolorbox}
        
    \begin{tcolorbox}[breakable, size=fbox, boxrule=1pt, pad at break*=1mm,colback=cellbackground, colframe=cellborder]
\prompt{In}{incolor}{13}{\boxspacing}
\begin{Verbatim}[commandchars=\\\{\}]
\PY{c+c1}{\PYZsh{} Valor da média}
\PY{n}{media\PYZus{}y} \PY{o}{=} \PY{p}{(}\PY{n}{freq\PYZus{}y}\PY{p}{[}\PY{l+s+s1}{\PYZsq{}}\PY{l+s+s1}{Y}\PY{l+s+s1}{\PYZsq{}}\PY{p}{]} \PY{o}{*} \PY{n}{freq\PYZus{}y}\PY{p}{[}\PY{l+s+s1}{\PYZsq{}}\PY{l+s+s1}{fr\PYZus{}y}\PY{l+s+s1}{\PYZsq{}}\PY{p}{]}\PY{p}{)}\PY{o}{.}\PY{n}{sum}\PY{p}{(}\PY{p}{)}
\PY{n+nb}{print}\PY{p}{(}\PY{l+s+sa}{f}\PY{l+s+s1}{\PYZsq{}}\PY{l+s+s1}{A media da variável Y é: }\PY{l+s+si}{\PYZob{}}\PY{n}{media\PYZus{}y}\PY{l+s+si}{:}\PY{l+s+s1}{.2f}\PY{l+s+si}{\PYZcb{}}\PY{l+s+s1}{\PYZsq{}}\PY{p}{)}

\PY{n}{var\PYZus{}y} \PY{o}{=} \PY{p}{(}\PY{n}{freq\PYZus{}y}\PY{p}{[}\PY{l+s+s1}{\PYZsq{}}\PY{l+s+s1}{fi\PYZus{}y}\PY{l+s+s1}{\PYZsq{}}\PY{p}{]}\PY{o}{*}\PY{p}{(}\PY{p}{(}\PY{n}{freq\PYZus{}y}\PY{p}{[}\PY{l+s+s1}{\PYZsq{}}\PY{l+s+s1}{Y}\PY{l+s+s1}{\PYZsq{}}\PY{p}{]}\PY{o}{\PYZhy{}}\PY{n}{media\PYZus{}y}\PY{p}{)}\PY{o}{*}\PY{o}{*}\PY{l+m+mi}{2}\PY{p}{)}\PY{o}{/}\PY{n}{n}\PY{p}{)}\PY{o}{.}\PY{n}{sum}\PY{p}{(}\PY{p}{)}
\PY{n+nb}{print}\PY{p}{(}\PY{l+s+sa}{f}\PY{l+s+s1}{\PYZsq{}}\PY{l+s+s1}{A variância de Y é: }\PY{l+s+si}{\PYZob{}}\PY{n}{var\PYZus{}y}\PY{l+s+si}{:}\PY{l+s+s1}{.2f}\PY{l+s+si}{\PYZcb{}}\PY{l+s+s1}{\PYZsq{}}\PY{p}{)}

\PY{n}{std\PYZus{}y} \PY{o}{=} \PY{n}{np}\PY{o}{.}\PY{n}{sqrt}\PY{p}{(}\PY{n}{var\PYZus{}y}\PY{p}{)}
\PY{n+nb}{print}\PY{p}{(}\PY{l+s+sa}{f}\PY{l+s+s1}{\PYZsq{}}\PY{l+s+s1}{O desvio padrão de Y é: }\PY{l+s+si}{\PYZob{}}\PY{n}{std\PYZus{}y}\PY{l+s+si}{:}\PY{l+s+s1}{.2f}\PY{l+s+si}{\PYZcb{}}\PY{l+s+s1}{\PYZsq{}}\PY{p}{)}

\PY{n}{cv\PYZus{}y} \PY{o}{=} \PY{n}{std\PYZus{}y} \PY{o}{/} \PY{n}{media\PYZus{}y}
\PY{n+nb}{print}\PY{p}{(}\PY{l+s+sa}{f}\PY{l+s+s1}{\PYZsq{}}\PY{l+s+s1}{o Coeficiente de Y é: }\PY{l+s+si}{\PYZob{}}\PY{n}{cv\PYZus{}y}\PY{l+s+si}{:}\PY{l+s+s1}{.2f}\PY{l+s+si}{\PYZcb{}}\PY{l+s+s1}{\PYZsq{}}\PY{p}{)}
\end{Verbatim}
\end{tcolorbox}

    \begin{Verbatim}[commandchars=\\\{\}]
A media da variável Y é: 107.50
A variância de Y é: 68.65
O desvio padrão de Y é: 8.29
o Coeficiente de Y é: 0.08
    \end{Verbatim}

    \subsubsection{Para a variável z}\label{para-a-variuxe1vel-z}

    \begin{tcolorbox}[breakable, size=fbox, boxrule=1pt, pad at break*=1mm,colback=cellbackground, colframe=cellborder]
\prompt{In}{incolor}{14}{\boxspacing}
\begin{Verbatim}[commandchars=\\\{\}]
\PY{n}{freq\PYZus{}z} \PY{o}{=} \PY{p}{(}
    \PY{n}{df}\PY{p}{[}\PY{l+s+s2}{\PYZdq{}}\PY{l+s+s2}{z}\PY{l+s+s2}{\PYZdq{}}\PY{p}{]}
    \PY{o}{.}\PY{n}{value\PYZus{}counts}\PY{p}{(}\PY{p}{)}
    \PY{o}{.}\PY{n}{sort\PYZus{}index}\PY{p}{(}\PY{p}{)}
    \PY{o}{.}\PY{n}{reset\PYZus{}index}\PY{p}{(}\PY{p}{)}
\PY{p}{)}

\PY{n}{freq\PYZus{}z}\PY{o}{.}\PY{n}{columns} \PY{o}{=} \PY{p}{[}\PY{l+s+s2}{\PYZdq{}}\PY{l+s+s2}{Z}\PY{l+s+s2}{\PYZdq{}}\PY{p}{,} \PY{l+s+s2}{\PYZdq{}}\PY{l+s+s2}{fi\PYZus{}z}\PY{l+s+s2}{\PYZdq{}}\PY{p}{]}

\PY{n}{freq\PYZus{}z}\PY{p}{[}\PY{l+s+s1}{\PYZsq{}}\PY{l+s+s1}{Fi\PYZus{}z}\PY{l+s+s1}{\PYZsq{}}\PY{p}{]} \PY{o}{=} \PY{n}{freq\PYZus{}z}\PY{p}{[}\PY{l+s+s1}{\PYZsq{}}\PY{l+s+s1}{fi\PYZus{}z}\PY{l+s+s1}{\PYZsq{}}\PY{p}{]}\PY{o}{.}\PY{n}{cumsum}\PY{p}{(}\PY{p}{)}
\PY{n}{freq\PYZus{}z}\PY{p}{[}\PY{l+s+s1}{\PYZsq{}}\PY{l+s+s1}{fr\PYZus{}z}\PY{l+s+s1}{\PYZsq{}}\PY{p}{]} \PY{o}{=} \PY{n}{freq\PYZus{}z}\PY{p}{[}\PY{l+s+s1}{\PYZsq{}}\PY{l+s+s1}{fi\PYZus{}z}\PY{l+s+s1}{\PYZsq{}}\PY{p}{]} \PY{o}{/} \PY{n}{n}
\PY{n}{freq\PYZus{}z}\PY{p}{[}\PY{l+s+s1}{\PYZsq{}}\PY{l+s+s1}{Frac\PYZus{}z}\PY{l+s+s1}{\PYZsq{}}\PY{p}{]} \PY{o}{=} \PY{n}{freq\PYZus{}z}\PY{p}{[}\PY{l+s+s1}{\PYZsq{}}\PY{l+s+s1}{fr\PYZus{}z}\PY{l+s+s1}{\PYZsq{}}\PY{p}{]}\PY{o}{.}\PY{n}{cumsum}\PY{p}{(}\PY{p}{)}

\PY{n}{freq\PYZus{}z}
\end{Verbatim}
\end{tcolorbox}

            \begin{tcolorbox}[breakable, size=fbox, boxrule=.5pt, pad at break*=1mm, opacityfill=0]
\prompt{Out}{outcolor}{14}{\boxspacing}
\begin{Verbatim}[commandchars=\\\{\}]
     Z  fi\_z  Fi\_z  fr\_z  Frac\_z
0   60     1     1  0.05    0.05
1   70     2     3  0.10    0.15
2   80     5     8  0.25    0.40
3   90     9    17  0.45    0.85
4  100     3    20  0.15    1.00
\end{Verbatim}
\end{tcolorbox}
        
    \begin{tcolorbox}[breakable, size=fbox, boxrule=1pt, pad at break*=1mm,colback=cellbackground, colframe=cellborder]
\prompt{In}{incolor}{15}{\boxspacing}
\begin{Verbatim}[commandchars=\\\{\}]
\PY{c+c1}{\PYZsh{} Valor da média}
\PY{n}{media\PYZus{}z} \PY{o}{=} \PY{p}{(}\PY{n}{freq\PYZus{}z}\PY{p}{[}\PY{l+s+s1}{\PYZsq{}}\PY{l+s+s1}{Z}\PY{l+s+s1}{\PYZsq{}}\PY{p}{]} \PY{o}{*} \PY{n}{freq\PYZus{}z}\PY{p}{[}\PY{l+s+s1}{\PYZsq{}}\PY{l+s+s1}{fr\PYZus{}z}\PY{l+s+s1}{\PYZsq{}}\PY{p}{]}\PY{p}{)}\PY{o}{.}\PY{n}{sum}\PY{p}{(}\PY{p}{)}
\PY{n+nb}{print}\PY{p}{(}\PY{l+s+sa}{f}\PY{l+s+s1}{\PYZsq{}}\PY{l+s+s1}{A media da variável Z é: }\PY{l+s+si}{\PYZob{}}\PY{n}{media\PYZus{}z}\PY{l+s+si}{:}\PY{l+s+s1}{.2f}\PY{l+s+si}{\PYZcb{}}\PY{l+s+s1}{\PYZsq{}}\PY{p}{)}

\PY{n}{var\PYZus{}z} \PY{o}{=} \PY{p}{(}\PY{p}{(}\PY{n}{freq\PYZus{}z}\PY{p}{[}\PY{l+s+s1}{\PYZsq{}}\PY{l+s+s1}{fi\PYZus{}z}\PY{l+s+s1}{\PYZsq{}}\PY{p}{]}\PY{o}{*}\PY{p}{(}\PY{n}{freq\PYZus{}z}\PY{p}{[}\PY{l+s+s1}{\PYZsq{}}\PY{l+s+s1}{Z}\PY{l+s+s1}{\PYZsq{}}\PY{p}{]}\PY{o}{\PYZhy{}}\PY{n}{media\PYZus{}z}\PY{p}{)}\PY{o}{*}\PY{o}{*}\PY{l+m+mi}{2}\PY{p}{)}\PY{o}{/}\PY{n}{n}\PY{p}{)}\PY{o}{.}\PY{n}{sum}\PY{p}{(}\PY{p}{)}
\PY{n+nb}{print}\PY{p}{(}\PY{l+s+sa}{f}\PY{l+s+s1}{\PYZsq{}}\PY{l+s+s1}{A variância de Z é: }\PY{l+s+si}{\PYZob{}}\PY{n}{var\PYZus{}z}\PY{l+s+si}{:}\PY{l+s+s1}{.2f}\PY{l+s+si}{\PYZcb{}}\PY{l+s+s1}{\PYZsq{}}\PY{p}{)}

\PY{n}{std\PYZus{}z} \PY{o}{=} \PY{n}{np}\PY{o}{.}\PY{n}{sqrt}\PY{p}{(}\PY{n}{var\PYZus{}z}\PY{p}{)}
\PY{n+nb}{print}\PY{p}{(}\PY{l+s+sa}{f}\PY{l+s+s1}{\PYZsq{}}\PY{l+s+s1}{O desvio padrão de Z é: }\PY{l+s+si}{\PYZob{}}\PY{n}{std\PYZus{}z}\PY{l+s+si}{:}\PY{l+s+s1}{.2f}\PY{l+s+si}{\PYZcb{}}\PY{l+s+s1}{\PYZsq{}}\PY{p}{)}

\PY{n}{cv\PYZus{}z} \PY{o}{=} \PY{n}{std\PYZus{}z} \PY{o}{/} \PY{n}{media\PYZus{}z}
\PY{n+nb}{print}\PY{p}{(}\PY{l+s+sa}{f}\PY{l+s+s1}{\PYZsq{}}\PY{l+s+s1}{o Coeficiente de Z é: }\PY{l+s+si}{\PYZob{}}\PY{n}{cv\PYZus{}z}\PY{l+s+si}{:}\PY{l+s+s1}{.2f}\PY{l+s+si}{\PYZcb{}}\PY{l+s+s1}{\PYZsq{}}\PY{p}{)}
\end{Verbatim}
\end{tcolorbox}

    \begin{Verbatim}[commandchars=\\\{\}]
A media da variável Z é: 85.50
A variância de Z é: 104.75
O desvio padrão de Z é: 10.23
o Coeficiente de Z é: 0.12
    \end{Verbatim}

    \subsection{Quantis de Y e Z}\label{quantis-de-y-e-z}

    \begin{tcolorbox}[breakable, size=fbox, boxrule=1pt, pad at break*=1mm,colback=cellbackground, colframe=cellborder]
\prompt{In}{incolor}{16}{\boxspacing}
\begin{Verbatim}[commandchars=\\\{\}]
\PY{n}{y\PYZus{}ord} \PY{o}{=} \PY{n}{np}\PY{o}{.}\PY{n}{array}\PY{p}{(}\PY{n+nb}{sorted}\PY{p}{(}\PY{n}{y}\PY{p}{)}\PY{p}{)}
\PY{n}{md} \PY{o}{=} \PY{p}{(}\PY{n}{y\PYZus{}ord}\PY{p}{[}\PY{l+m+mi}{9}\PY{p}{]}\PY{o}{+}\PY{n}{y\PYZus{}ord}\PY{p}{[}\PY{l+m+mi}{10}\PY{p}{]}\PY{p}{)}\PY{o}{/}\PY{l+m+mi}{2}
\PY{n+nb}{print}\PY{p}{(}\PY{n}{y\PYZus{}ord}\PY{p}{)}
\PY{n+nb}{print}\PY{p}{(}\PY{l+s+sa}{f}\PY{l+s+s1}{\PYZsq{}}\PY{l+s+s1}{A mediana é: }\PY{l+s+si}{\PYZob{}}\PY{n}{md}\PY{l+s+si}{:}\PY{l+s+s1}{.2f}\PY{l+s+si}{\PYZcb{}}\PY{l+s+s1}{\PYZsq{}}\PY{p}{)}
\PY{n}{q1} \PY{o}{=} \PY{n}{y\PYZus{}ord}\PY{p}{[}\PY{l+m+mi}{4}\PY{p}{]}
\PY{n+nb}{print}\PY{p}{(}\PY{l+s+sa}{f}\PY{l+s+s1}{\PYZsq{}}\PY{l+s+s1}{o valor de Q1 é: }\PY{l+s+si}{\PYZob{}}\PY{n}{q1}\PY{l+s+si}{:}\PY{l+s+s1}{.2f}\PY{l+s+si}{\PYZcb{}}\PY{l+s+s1}{\PYZsq{}}\PY{p}{)}
\PY{n}{q3} \PY{o}{=} \PY{p}{(}\PY{n}{y\PYZus{}ord}\PY{p}{[}\PY{l+m+mi}{14}\PY{p}{]}\PY{p}{)}
\PY{n+nb}{print}\PY{p}{(}\PY{l+s+sa}{f}\PY{l+s+s1}{\PYZsq{}}\PY{l+s+s1}{o valor de Q3 é: }\PY{l+s+si}{\PYZob{}}\PY{n}{q3}\PY{l+s+si}{:}\PY{l+s+s1}{.2f}\PY{l+s+si}{\PYZcb{}}\PY{l+s+s1}{\PYZsq{}}\PY{p}{)}
\end{Verbatim}
\end{tcolorbox}

    \begin{Verbatim}[commandchars=\\\{\}]
[ 92  96  98 100 100 101 104 105 106 106 108 109 110 112 112 113 116 117
 118 127]
A mediana é: 107.00
o valor de Q1 é: 100.00
o valor de Q3 é: 112.00
    \end{Verbatim}

    \begin{tcolorbox}[breakable, size=fbox, boxrule=1pt, pad at break*=1mm,colback=cellbackground, colframe=cellborder]
\prompt{In}{incolor}{17}{\boxspacing}
\begin{Verbatim}[commandchars=\\\{\}]
\PY{n}{plt}\PY{o}{.}\PY{n}{boxplot}\PY{p}{(}\PY{n}{y\PYZus{}ord}\PY{p}{,} \PY{n}{orientation}\PY{o}{=}\PY{l+s+s1}{\PYZsq{}}\PY{l+s+s1}{horizontal}\PY{l+s+s1}{\PYZsq{}}\PY{p}{,} \PY{n}{showmeans}\PY{o}{=}\PY{k+kc}{True}\PY{p}{)}\PY{p}{;}
\PY{n}{plt}\PY{o}{.}\PY{n}{show}\PY{p}{(}\PY{p}{)}
\end{Verbatim}
\end{tcolorbox}

    \begin{center}
    \adjustimage{max size={0.9\linewidth}{0.9\paperheight}}{Exercicios_estatistica_rev_verificada_files/Exercicios_estatistica_rev_verificada_28_0.png}
    \end{center}
    { \hspace*{\fill} \\}
    
    \begin{tcolorbox}[breakable, size=fbox, boxrule=1pt, pad at break*=1mm,colback=cellbackground, colframe=cellborder]
\prompt{In}{incolor}{18}{\boxspacing}
\begin{Verbatim}[commandchars=\\\{\}]
\PY{n}{z\PYZus{}ord} \PY{o}{=} \PY{n}{np}\PY{o}{.}\PY{n}{array}\PY{p}{(}\PY{n+nb}{sorted}\PY{p}{(}\PY{n}{z}\PY{p}{)}\PY{p}{)}

\PY{n}{md\PYZus{}z} \PY{o}{=} \PY{p}{(}\PY{n}{z\PYZus{}ord}\PY{p}{[}\PY{l+m+mi}{9}\PY{p}{]} \PY{o}{+} \PY{n}{z\PYZus{}ord}\PY{p}{[}\PY{l+m+mi}{10}\PY{p}{]}\PY{p}{)} \PY{o}{/} \PY{l+m+mi}{2}
\PY{n}{q1\PYZus{}z} \PY{o}{=} \PY{n}{z\PYZus{}ord}\PY{p}{[}\PY{l+m+mi}{4}\PY{p}{]}
\PY{n}{q3\PYZus{}z} \PY{o}{=} \PY{n}{z\PYZus{}ord}\PY{p}{[}\PY{l+m+mi}{14}\PY{p}{]}

\PY{n+nb}{print}\PY{p}{(}\PY{n}{z\PYZus{}ord}\PY{p}{)}
\PY{n+nb}{print}\PY{p}{(}\PY{l+s+sa}{f}\PY{l+s+s1}{\PYZsq{}}\PY{l+s+s1}{A mediana é: }\PY{l+s+si}{\PYZob{}}\PY{n}{md\PYZus{}z}\PY{l+s+si}{:}\PY{l+s+s1}{.2f}\PY{l+s+si}{\PYZcb{}}\PY{l+s+s1}{\PYZsq{}}\PY{p}{)}
\PY{n+nb}{print}\PY{p}{(}\PY{l+s+sa}{f}\PY{l+s+s1}{\PYZsq{}}\PY{l+s+s1}{o valor de Q1 é: }\PY{l+s+si}{\PYZob{}}\PY{n}{q1\PYZus{}z}\PY{l+s+si}{:}\PY{l+s+s1}{.2f}\PY{l+s+si}{\PYZcb{}}\PY{l+s+s1}{\PYZsq{}}\PY{p}{)}
\PY{n+nb}{print}\PY{p}{(}\PY{l+s+sa}{f}\PY{l+s+s1}{\PYZsq{}}\PY{l+s+s1}{o valor de Q3 é: }\PY{l+s+si}{\PYZob{}}\PY{n}{q3\PYZus{}z}\PY{l+s+si}{:}\PY{l+s+s1}{.2f}\PY{l+s+si}{\PYZcb{}}\PY{l+s+s1}{\PYZsq{}}\PY{p}{)}
\end{Verbatim}
\end{tcolorbox}

    \begin{Verbatim}[commandchars=\\\{\}]
[ 60  70  70  80  80  80  80  80  90  90  90  90  90  90  90  90  90 100
 100 100]
A mediana é: 90.00
o valor de Q1 é: 80.00
o valor de Q3 é: 90.00
    \end{Verbatim}

    \begin{tcolorbox}[breakable, size=fbox, boxrule=1pt, pad at break*=1mm,colback=cellbackground, colframe=cellborder]
\prompt{In}{incolor}{19}{\boxspacing}
\begin{Verbatim}[commandchars=\\\{\}]
\PY{n}{plt}\PY{o}{.}\PY{n}{boxplot}\PY{p}{(}\PY{n}{z\PYZus{}ord}\PY{p}{,} \PY{n}{orientation}\PY{o}{=}\PY{l+s+s1}{\PYZsq{}}\PY{l+s+s1}{horizontal}\PY{l+s+s1}{\PYZsq{}}\PY{p}{,} \PY{n}{showmeans}\PY{o}{=}\PY{k+kc}{True}\PY{p}{)}\PY{p}{;}
\PY{n}{plt}\PY{o}{.}\PY{n}{show}\PY{p}{(}\PY{p}{)}
\end{Verbatim}
\end{tcolorbox}

    \begin{center}
    \adjustimage{max size={0.9\linewidth}{0.9\paperheight}}{Exercicios_estatistica_rev_verificada_files/Exercicios_estatistica_rev_verificada_30_0.png}
    \end{center}
    { \hspace*{\fill} \\}
    
    4 -- Após um tratamento, seis operários submeteram-se a um teste e, mais
tarde, mediu-se a produtividade de cada um deles. A partir dos
resultados apresentados na tabela abaixo, calcule o coeficiente de
correlação entre a nota do teste e a produtividade. O que podemos aferir
desse resultado? Calcule também a assimetria da amostra teste e da
amostra produtividade e compare-as.

{\def\LTcaptype{table} % do not increment counter
\begin{longtable}[]{@{}rrrrrrr@{}}
\toprule\noalign{}
Operário & 1 & 2 & 3 & 4 & 5 & 6 \\
\midrule\noalign{}
\endhead
\bottomrule\noalign{}
\endlastfoot
Teste & 9 & 17 & 20 & 19 & 20 & 23 \\
Produtividade & 22 & 34 & 29 & 33 & 42 & 32 \\
\end{longtable}
}

    5-- Procurando quantificar os efeitos da escassez de sono sobre a
capacidade de resolução de problemas simples, um agente tomou ao acaso
10 sujeitos e os submeteu a experimentação. Deixou-os sem dormir por
diferentes números de horas, após o que solicitou que os mesmos
resolvessem os itens ``contas de adicionar'' de um teste. Obteve, assim,
os seguintes dados:

{\def\LTcaptype{table} % do not increment counter
\begin{longtable}[]{@{}rr@{}}
\toprule\noalign{}
Nº de erros --- Y & Horas sem dormir --- X \\
\midrule\noalign{}
\endhead
\bottomrule\noalign{}
\endlastfoot
8 & 8 \\
6 & 8 \\
6 & 12 \\
10 & 12 \\
8 & 16 \\
14 & 16 \\
14 & 20 \\
12 & 20 \\
16 & 24 \\
12 & 24 \\
\end{longtable}
}

Calcule o coeficiente de correlação. Pode-se concluir que horas sem
dormir afeta a capacidade de um sujeito durante um teste?

    6 -- Em um exame de colesterol, o grau médio de um grupo ``A'' de 150
pessoas foi de 214 mg/dl e um desvio-padrão de 22 mg/dl. Em outro grupo
``B'', entretanto, o grau médio em 150 pessoas foi de 201 mg/dl e um
desvio-padrão de 21 mg/dl. Em que grupo foi maior a dispersão?

{\def\LTcaptype{table} % do not increment counter
\begin{longtable}[]{@{}rr@{}}
\toprule\noalign{}
Renda familiar (X) & Gasto com alimentação (Y) \\
\midrule\noalign{}
\endhead
\bottomrule\noalign{}
\endlastfoot
3 & 1,5 \\
5 & 2,0 \\
10 & 6,0 \\
10 & 7,0 \\
20 & 10,0 \\
20 & 12,0 \\
20 & 15,0 \\
30 & 8,0 \\
40 & 10,0 \\
50 & 20,0 \\
60 & 20,0 \\
70 & 25,0 \\
70 & 30,0 \\
80 & 25,0 \\
100 & 40,0 \\
100 & 35,0 \\
100 & 40,0 \\
120 & 30,0 \\
120 & 40,0 \\
140 & 40,0 \\
150 & 50,0 \\
180 & 40,0 \\
180 & 50,0 \\
200 & 60,0 \\
200 & 50,0 \\
\end{longtable}
}

    \begin{itemize}
\tightlist
\item
  Calcule a distribuição de frequências das duas variáveis aleatórias.\\
\item
  \st{Calcule a distribuição de frequências conjunta das duas variáveis
  aleatórias.}\\
\item
  Calcule o coeficiente de variação de cada uma das duas variáveis
  aleatórias e interprete-os.
\item
  Calcule o coeficiente de correlação entre a renda familiar e o gasto
  com alimentação desse conjunto de indivíduos. Interprete-o.
\item
  Calcule os quartis das duas distribuições de frequências. Monte o
  diagrama de Box-Plot.
\end{itemize}

    7 - renda familiar e gastos com alimentação.

    \begin{tcolorbox}[breakable, size=fbox, boxrule=1pt, pad at break*=1mm,colback=cellbackground, colframe=cellborder]
\prompt{In}{incolor}{20}{\boxspacing}
\begin{Verbatim}[commandchars=\\\{\}]
\PY{n}{renda\PYZus{}familiar} \PY{o}{=} \PY{p}{[}
    \PY{l+m+mi}{3}\PY{p}{,} \PY{l+m+mi}{5}\PY{p}{,} \PY{l+m+mi}{10}\PY{p}{,} \PY{l+m+mi}{10}\PY{p}{,} \PY{l+m+mi}{20}\PY{p}{,} \PY{l+m+mi}{20}\PY{p}{,}
    \PY{l+m+mi}{20}\PY{p}{,} \PY{l+m+mi}{30}\PY{p}{,} \PY{l+m+mi}{40}\PY{p}{,} \PY{l+m+mi}{50}\PY{p}{,} \PY{l+m+mi}{60}\PY{p}{,} \PY{l+m+mi}{70}\PY{p}{,}
    \PY{l+m+mi}{70}\PY{p}{,} \PY{l+m+mi}{80}\PY{p}{,} \PY{l+m+mi}{100}\PY{p}{,} \PY{l+m+mi}{100}\PY{p}{,} \PY{l+m+mi}{100}\PY{p}{,}
    \PY{l+m+mi}{120}\PY{p}{,} \PY{l+m+mi}{120}\PY{p}{,} \PY{l+m+mi}{140}\PY{p}{,} \PY{l+m+mi}{150}\PY{p}{,} \PY{l+m+mi}{180}\PY{p}{,}
    \PY{l+m+mi}{180}\PY{p}{,} \PY{l+m+mi}{200}\PY{p}{,} \PY{l+m+mi}{200}
\PY{p}{]}

\PY{n}{n\PYZus{}renda} \PY{o}{=} \PY{n+nb}{len}\PY{p}{(}\PY{n}{renda\PYZus{}familiar}\PY{p}{)}

\PY{n}{dados\PYZus{}renda} \PY{o}{=} \PY{n}{pd}\PY{o}{.}\PY{n}{DataFrame}\PY{p}{(}\PY{p}{\PYZob{}}
    \PY{l+s+s2}{\PYZdq{}}\PY{l+s+s2}{Renda Familiar}\PY{l+s+s2}{\PYZdq{}}\PY{p}{:} \PY{n}{renda\PYZus{}familiar}
\PY{p}{\PYZcb{}}\PY{p}{)}

\PY{n}{freq\PYZus{}renda} \PY{o}{=} \PY{p}{(}
    \PY{n}{dados\PYZus{}renda}\PY{p}{[}\PY{l+s+s2}{\PYZdq{}}\PY{l+s+s2}{Renda Familiar}\PY{l+s+s2}{\PYZdq{}}\PY{p}{]}
    \PY{o}{.}\PY{n}{value\PYZus{}counts}\PY{p}{(}\PY{p}{)}
    \PY{o}{.}\PY{n}{sort\PYZus{}index}\PY{p}{(}\PY{p}{)}
    \PY{o}{.}\PY{n}{reset\PYZus{}index}\PY{p}{(}\PY{p}{)}
\PY{p}{)}

\PY{n}{freq\PYZus{}renda}\PY{o}{.}\PY{n}{columns} \PY{o}{=} \PY{p}{[}\PY{l+s+s2}{\PYZdq{}}\PY{l+s+s2}{Renda}\PY{l+s+s2}{\PYZdq{}}\PY{p}{,} \PY{l+s+s2}{\PYZdq{}}\PY{l+s+s2}{fi\PYZus{}renda}\PY{l+s+s2}{\PYZdq{}}\PY{p}{]}

\PY{n}{freq\PYZus{}renda}\PY{p}{[}\PY{l+s+s2}{\PYZdq{}}\PY{l+s+s2}{Fi\PYZus{}renda}\PY{l+s+s2}{\PYZdq{}}\PY{p}{]} \PY{o}{=} \PY{n}{freq\PYZus{}renda}\PY{p}{[}\PY{l+s+s2}{\PYZdq{}}\PY{l+s+s2}{fi\PYZus{}renda}\PY{l+s+s2}{\PYZdq{}}\PY{p}{]}\PY{o}{.}\PY{n}{cumsum}\PY{p}{(}\PY{p}{)}
\PY{n}{freq\PYZus{}renda}\PY{p}{[}\PY{l+s+s2}{\PYZdq{}}\PY{l+s+s2}{fr\PYZus{}renda}\PY{l+s+s2}{\PYZdq{}}\PY{p}{]} \PY{o}{=} \PY{n}{freq\PYZus{}renda}\PY{p}{[}\PY{l+s+s2}{\PYZdq{}}\PY{l+s+s2}{fi\PYZus{}renda}\PY{l+s+s2}{\PYZdq{}}\PY{p}{]} \PY{o}{/} \PY{n}{n\PYZus{}renda}
\PY{n}{freq\PYZus{}renda}\PY{p}{[}\PY{l+s+s2}{\PYZdq{}}\PY{l+s+s2}{Fr\PYZus{}renda}\PY{l+s+s2}{\PYZdq{}}\PY{p}{]} \PY{o}{=} \PY{n}{freq\PYZus{}renda}\PY{p}{[}\PY{l+s+s2}{\PYZdq{}}\PY{l+s+s2}{fr\PYZus{}renda}\PY{l+s+s2}{\PYZdq{}}\PY{p}{]}\PY{o}{.}\PY{n}{cumsum}\PY{p}{(}\PY{p}{)}

\PY{n}{freq\PYZus{}renda}
\end{Verbatim}
\end{tcolorbox}

            \begin{tcolorbox}[breakable, size=fbox, boxrule=.5pt, pad at break*=1mm, opacityfill=0]
\prompt{Out}{outcolor}{20}{\boxspacing}
\begin{Verbatim}[commandchars=\\\{\}]
    Renda  fi\_renda  Fi\_renda  fr\_renda  Fr\_renda
0       3         1         1      0.04      0.04
1       5         1         2      0.04      0.08
2      10         2         4      0.08      0.16
3      20         3         7      0.12      0.28
4      30         1         8      0.04      0.32
5      40         1         9      0.04      0.36
6      50         1        10      0.04      0.40
7      60         1        11      0.04      0.44
8      70         2        13      0.08      0.52
9      80         1        14      0.04      0.56
10    100         3        17      0.12      0.68
11    120         2        19      0.08      0.76
12    140         1        20      0.04      0.80
13    150         1        21      0.04      0.84
14    180         2        23      0.08      0.92
15    200         2        25      0.08      1.00
\end{Verbatim}
\end{tcolorbox}
        
    \subsubsection{Estatísticas.}\label{estatuxedsticas.}

a média é:

\[ \bar{x} = \sum(fr_i \cdot x_i) \]

    a variância populacional é:

\[
\sigma^2 =
\frac{\sum f_i (x_i - \bar{x})^2}{n}
\]

e o desvio padrão:

\[
\sigma = \sqrt{\sigma^2}
\]

e a assimetria pelo terceiro momento padronizado:

\[
g =
\frac{\sum f_i (x_i - \bar{x})^3}
{n\,\sigma^3}
\]

    \begin{tcolorbox}[breakable, size=fbox, boxrule=1pt, pad at break*=1mm,colback=cellbackground, colframe=cellborder]
\prompt{In}{incolor}{21}{\boxspacing}
\begin{Verbatim}[commandchars=\\\{\}]
\PY{c+c1}{\PYZsh{} média}

\PY{n}{media\PYZus{}renda} \PY{o}{=} \PY{p}{(}
    \PY{n}{freq\PYZus{}renda}\PY{p}{[}\PY{l+s+s2}{\PYZdq{}}\PY{l+s+s2}{Renda}\PY{l+s+s2}{\PYZdq{}}\PY{p}{]} \PY{o}{*} \PY{n}{freq\PYZus{}renda}\PY{p}{[}\PY{l+s+s2}{\PYZdq{}}\PY{l+s+s2}{fr\PYZus{}renda}\PY{l+s+s2}{\PYZdq{}}\PY{p}{]}
\PY{p}{)}\PY{o}{.}\PY{n}{sum}\PY{p}{(}\PY{p}{)}

\PY{n}{var\PYZus{}renda} \PY{o}{=} \PY{p}{(}
    \PY{n}{freq\PYZus{}renda}\PY{p}{[}\PY{l+s+s2}{\PYZdq{}}\PY{l+s+s2}{fi\PYZus{}renda}\PY{l+s+s2}{\PYZdq{}}\PY{p}{]}
    \PY{o}{*} \PY{p}{(}\PY{n}{freq\PYZus{}renda}\PY{p}{[}\PY{l+s+s2}{\PYZdq{}}\PY{l+s+s2}{Renda}\PY{l+s+s2}{\PYZdq{}}\PY{p}{]} \PY{o}{\PYZhy{}} \PY{n}{media\PYZus{}renda}\PY{p}{)}\PY{o}{*}\PY{o}{*}\PY{l+m+mi}{2}
\PY{p}{)}\PY{o}{.}\PY{n}{sum}\PY{p}{(}\PY{p}{)} \PY{o}{/} \PY{n}{n\PYZus{}renda}

\PY{n}{std\PYZus{}renda} \PY{o}{=} \PY{n}{np}\PY{o}{.}\PY{n}{sqrt}\PY{p}{(}\PY{n}{var\PYZus{}renda}\PY{p}{)}

\PY{n}{cv\PYZus{}renda} \PY{o}{=} \PY{n}{std\PYZus{}renda} \PY{o}{/} \PY{n}{media\PYZus{}renda} \PY{o}{*} \PY{l+m+mi}{100}

\PY{n}{g\PYZus{}renda} \PY{o}{=} \PY{p}{(}
    \PY{n}{freq\PYZus{}renda}\PY{p}{[}\PY{l+s+s2}{\PYZdq{}}\PY{l+s+s2}{fi\PYZus{}renda}\PY{l+s+s2}{\PYZdq{}}\PY{p}{]}
    \PY{o}{*} \PY{p}{(}\PY{n}{freq\PYZus{}renda}\PY{p}{[}\PY{l+s+s2}{\PYZdq{}}\PY{l+s+s2}{Renda}\PY{l+s+s2}{\PYZdq{}}\PY{p}{]} \PY{o}{\PYZhy{}} \PY{n}{media\PYZus{}renda}\PY{p}{)}\PY{o}{*}\PY{o}{*}\PY{l+m+mi}{3}
\PY{p}{)}\PY{o}{.}\PY{n}{sum}\PY{p}{(}\PY{p}{)} \PY{o}{/} \PY{p}{(}\PY{n}{n\PYZus{}renda} \PY{o}{*} \PY{n}{std\PYZus{}renda}\PY{o}{*}\PY{o}{*}\PY{l+m+mi}{3}\PY{p}{)}

\PY{n+nb}{print}\PY{p}{(}\PY{l+s+sa}{f}\PY{l+s+s1}{\PYZsq{}}\PY{l+s+s1}{O valor da média da renda é: }\PY{l+s+si}{\PYZob{}}\PY{n}{media\PYZus{}renda}\PY{l+s+si}{:}\PY{l+s+s1}{.2f}\PY{l+s+si}{\PYZcb{}}\PY{l+s+s1}{\PYZsq{}}\PY{p}{)}
\PY{n+nb}{print}\PY{p}{(}\PY{l+s+sa}{f}\PY{l+s+s1}{\PYZsq{}}\PY{l+s+s1}{O valor da variância da renda é: }\PY{l+s+si}{\PYZob{}}\PY{n}{var\PYZus{}renda}\PY{l+s+si}{:}\PY{l+s+s1}{.2f}\PY{l+s+si}{\PYZcb{}}\PY{l+s+s1}{\PYZsq{}}\PY{p}{)}
\PY{n+nb}{print}\PY{p}{(}\PY{l+s+sa}{f}\PY{l+s+s1}{\PYZsq{}}\PY{l+s+s1}{O valor do desvio padrão da renda é: }\PY{l+s+si}{\PYZob{}}\PY{n}{std\PYZus{}renda}\PY{l+s+si}{:}\PY{l+s+s1}{.2f}\PY{l+s+si}{\PYZcb{}}\PY{l+s+s1}{\PYZsq{}}\PY{p}{)}
\PY{n+nb}{print}\PY{p}{(}\PY{l+s+sa}{f}\PY{l+s+s1}{\PYZsq{}}\PY{l+s+s1}{O valor do coeficiente de variação é: }\PY{l+s+si}{\PYZob{}}\PY{n}{cv\PYZus{}renda}\PY{l+s+si}{:}\PY{l+s+s1}{.2f}\PY{l+s+si}{\PYZcb{}}\PY{l+s+s1}{\PYZpc{}}\PY{l+s+s1}{\PYZsq{}}\PY{p}{)}
\PY{n+nb}{print}\PY{p}{(}\PY{l+s+sa}{f}\PY{l+s+s1}{\PYZsq{}}\PY{l+s+s1}{O valor da assimetria é: }\PY{l+s+si}{\PYZob{}}\PY{n}{g\PYZus{}renda}\PY{l+s+si}{:}\PY{l+s+s1}{.2f}\PY{l+s+si}{\PYZcb{}}\PY{l+s+s1}{\PYZsq{}}\PY{p}{)}
\end{Verbatim}
\end{tcolorbox}

    \begin{Verbatim}[commandchars=\\\{\}]
O valor da média da renda é: 83.12
O valor da variância da renda é: 3968.43
O valor do desvio padrão da renda é: 63.00
O valor do coeficiente de variação é: 75.79\%
O valor da assimetria é: 0.45
    \end{Verbatim}

    \emph{O coeficiente de variação (CV) mede a dispersão dos dados em
relação à média. Um CV baixo indica alta homogeneidade e estabilidade,
enquanto um CV alto aponta grande dispersão e instabilidade. O valor
padrão geral classifica variações até 10\% como baixas, e acima de 30\%
como muito altas.}

    \subsection{Quantis da renda
familiar.}\label{quantis-da-renda-familiar.}

    \begin{tcolorbox}[breakable, size=fbox, boxrule=1pt, pad at break*=1mm,colback=cellbackground, colframe=cellborder]
\prompt{In}{incolor}{22}{\boxspacing}
\begin{Verbatim}[commandchars=\\\{\}]
\PY{n}{renda\PYZus{}ord} \PY{o}{=} \PY{n+nb}{sorted}\PY{p}{(}\PY{n}{renda\PYZus{}familiar}\PY{p}{)}

\PY{n}{mediana\PYZus{}renda\PYZus{}familiar} \PY{o}{=} \PY{n}{renda\PYZus{}ord}\PY{p}{[}\PY{l+m+mi}{12}\PY{p}{]}
\PY{n}{q1\PYZus{}renda\PYZus{}familiar} \PY{o}{=} \PY{p}{(}\PY{n}{renda\PYZus{}ord}\PY{p}{[}\PY{l+m+mi}{5}\PY{p}{]} \PY{o}{+} \PY{n}{renda\PYZus{}ord}\PY{p}{[}\PY{l+m+mi}{6}\PY{p}{]}\PY{p}{)} \PY{o}{/} \PY{l+m+mi}{2}
\PY{n}{q3\PYZus{}renda\PYZus{}familiar} \PY{o}{=} \PY{p}{(}\PY{n}{renda\PYZus{}ord}\PY{p}{[}\PY{l+m+mi}{18}\PY{p}{]} \PY{o}{+} \PY{n}{renda\PYZus{}ord}\PY{p}{[}\PY{l+m+mi}{19}\PY{p}{]}\PY{p}{)} \PY{o}{/} \PY{l+m+mi}{2}

\PY{n+nb}{print}\PY{p}{(}\PY{n}{renda\PYZus{}ord}\PY{p}{)}
\PY{n+nb}{print}\PY{p}{(}\PY{l+s+sa}{f}\PY{l+s+s1}{\PYZsq{}}\PY{l+s+s1}{A mediana da renda é: }\PY{l+s+si}{\PYZob{}}\PY{n}{mediana\PYZus{}renda\PYZus{}familiar}\PY{l+s+si}{:}\PY{l+s+s1}{.2f}\PY{l+s+si}{\PYZcb{}}\PY{l+s+s1}{\PYZsq{}}\PY{p}{)}
\PY{n+nb}{print}\PY{p}{(}\PY{l+s+sa}{f}\PY{l+s+s1}{\PYZsq{}}\PY{l+s+s1}{O valor de Q1 da renda é: }\PY{l+s+si}{\PYZob{}}\PY{n}{q1\PYZus{}renda\PYZus{}familiar}\PY{l+s+si}{:}\PY{l+s+s1}{.2f}\PY{l+s+si}{\PYZcb{}}\PY{l+s+s1}{\PYZsq{}}\PY{p}{)}
\PY{n+nb}{print}\PY{p}{(}\PY{l+s+sa}{f}\PY{l+s+s1}{\PYZsq{}}\PY{l+s+s1}{O valor de Q3 da renda é: }\PY{l+s+si}{\PYZob{}}\PY{n}{q3\PYZus{}renda\PYZus{}familiar}\PY{l+s+si}{:}\PY{l+s+s1}{.2f}\PY{l+s+si}{\PYZcb{}}\PY{l+s+s1}{\PYZsq{}}\PY{p}{)}
\end{Verbatim}
\end{tcolorbox}

    \begin{Verbatim}[commandchars=\\\{\}]
[3, 5, 10, 10, 20, 20, 20, 30, 40, 50, 60, 70, 70, 80, 100, 100, 100, 120, 120,
140, 150, 180, 180, 200, 200]
A mediana da renda é: 70.00
O valor de Q1 da renda é: 20.00
O valor de Q3 da renda é: 130.00
    \end{Verbatim}

    Para o gráfico eu vou usar o recurso do Python.

    \begin{tcolorbox}[breakable, size=fbox, boxrule=1pt, pad at break*=1mm,colback=cellbackground, colframe=cellborder]
\prompt{In}{incolor}{23}{\boxspacing}
\begin{Verbatim}[commandchars=\\\{\}]
\PY{n}{plt}\PY{o}{.}\PY{n}{boxplot}\PY{p}{(}\PY{n}{renda\PYZus{}familiar}\PY{p}{,} \PY{n}{orientation}\PY{o}{=}\PY{l+s+s1}{\PYZsq{}}\PY{l+s+s1}{horizontal}\PY{l+s+s1}{\PYZsq{}}\PY{p}{,} \PY{n}{showmeans}\PY{o}{=}\PY{k+kc}{True}\PY{p}{)}\PY{p}{;}
\PY{n}{plt}\PY{o}{.}\PY{n}{show}\PY{p}{(}\PY{p}{)}
\end{Verbatim}
\end{tcolorbox}

    \begin{center}
    \adjustimage{max size={0.9\linewidth}{0.9\paperheight}}{Exercicios_estatistica_rev_verificada_files/Exercicios_estatistica_rev_verificada_44_0.png}
    \end{center}
    { \hspace*{\fill} \\}
    
    \subsection{Variável Y custos da
alimentação}\label{variuxe1vel-y-custos-da-alimentauxe7uxe3o}

    \begin{tcolorbox}[breakable, size=fbox, boxrule=1pt, pad at break*=1mm,colback=cellbackground, colframe=cellborder]
\prompt{In}{incolor}{24}{\boxspacing}
\begin{Verbatim}[commandchars=\\\{\}]
\PY{n}{aliment} \PY{o}{=} \PY{p}{[}
    \PY{l+m+mf}{1.5}\PY{p}{,} \PY{l+m+mi}{2}\PY{p}{,} \PY{l+m+mi}{6}\PY{p}{,} \PY{l+m+mi}{7}\PY{p}{,} \PY{l+m+mi}{10}\PY{p}{,} \PY{l+m+mi}{12}\PY{p}{,}
    \PY{l+m+mi}{15}\PY{p}{,} \PY{l+m+mi}{8}\PY{p}{,} \PY{l+m+mi}{10}\PY{p}{,} \PY{l+m+mi}{20}\PY{p}{,} \PY{l+m+mi}{20}\PY{p}{,} \PY{l+m+mi}{25}\PY{p}{,}
    \PY{l+m+mi}{30}\PY{p}{,} \PY{l+m+mi}{25}\PY{p}{,} \PY{l+m+mi}{40}\PY{p}{,} \PY{l+m+mi}{35}\PY{p}{,} \PY{l+m+mi}{40}\PY{p}{,}
    \PY{l+m+mi}{30}\PY{p}{,} \PY{l+m+mi}{40}\PY{p}{,} \PY{l+m+mi}{40}\PY{p}{,} \PY{l+m+mi}{50}\PY{p}{,} \PY{l+m+mi}{40}\PY{p}{,}
    \PY{l+m+mi}{50}\PY{p}{,} \PY{l+m+mi}{60}\PY{p}{,} \PY{l+m+mi}{50}
\PY{p}{]}

\PY{n}{n\PYZus{}aliment} \PY{o}{=} \PY{n+nb}{len}\PY{p}{(}\PY{n}{aliment}\PY{p}{)}

\PY{n}{dados\PYZus{}aliment} \PY{o}{=} \PY{n}{pd}\PY{o}{.}\PY{n}{DataFrame}\PY{p}{(}\PY{p}{\PYZob{}}
    \PY{l+s+s2}{\PYZdq{}}\PY{l+s+s2}{Gastos com Alimentação}\PY{l+s+s2}{\PYZdq{}}\PY{p}{:} \PY{n}{aliment}
\PY{p}{\PYZcb{}}\PY{p}{)}

\PY{n}{freq\PYZus{}aliment} \PY{o}{=} \PY{p}{(}
    \PY{n}{dados\PYZus{}aliment}\PY{p}{[}\PY{l+s+s2}{\PYZdq{}}\PY{l+s+s2}{Gastos com Alimentação}\PY{l+s+s2}{\PYZdq{}}\PY{p}{]}
    \PY{o}{.}\PY{n}{value\PYZus{}counts}\PY{p}{(}\PY{p}{)}
    \PY{o}{.}\PY{n}{sort\PYZus{}index}\PY{p}{(}\PY{p}{)}
    \PY{o}{.}\PY{n}{reset\PYZus{}index}\PY{p}{(}\PY{p}{)}
\PY{p}{)}

\PY{n}{freq\PYZus{}aliment}\PY{o}{.}\PY{n}{columns} \PY{o}{=} \PY{p}{[}\PY{l+s+s2}{\PYZdq{}}\PY{l+s+s2}{Gastos com Alimentação}\PY{l+s+s2}{\PYZdq{}}\PY{p}{,} \PY{l+s+s2}{\PYZdq{}}\PY{l+s+s2}{fi\PYZus{}aliment}\PY{l+s+s2}{\PYZdq{}}\PY{p}{]}

\PY{n}{freq\PYZus{}aliment}\PY{p}{[}\PY{l+s+s2}{\PYZdq{}}\PY{l+s+s2}{Fi\PYZus{}aliment}\PY{l+s+s2}{\PYZdq{}}\PY{p}{]} \PY{o}{=} \PY{n}{freq\PYZus{}aliment}\PY{p}{[}\PY{l+s+s2}{\PYZdq{}}\PY{l+s+s2}{fi\PYZus{}aliment}\PY{l+s+s2}{\PYZdq{}}\PY{p}{]}\PY{o}{.}\PY{n}{cumsum}\PY{p}{(}\PY{p}{)}
\PY{n}{freq\PYZus{}aliment}\PY{p}{[}\PY{l+s+s2}{\PYZdq{}}\PY{l+s+s2}{fr\PYZus{}aliment}\PY{l+s+s2}{\PYZdq{}}\PY{p}{]} \PY{o}{=} \PY{n}{freq\PYZus{}aliment}\PY{p}{[}\PY{l+s+s2}{\PYZdq{}}\PY{l+s+s2}{fi\PYZus{}aliment}\PY{l+s+s2}{\PYZdq{}}\PY{p}{]} \PY{o}{/} \PY{n}{n\PYZus{}aliment}
\PY{n}{freq\PYZus{}aliment}\PY{p}{[}\PY{l+s+s2}{\PYZdq{}}\PY{l+s+s2}{Fr\PYZus{}aliment}\PY{l+s+s2}{\PYZdq{}}\PY{p}{]} \PY{o}{=} \PY{n}{freq\PYZus{}aliment}\PY{p}{[}\PY{l+s+s2}{\PYZdq{}}\PY{l+s+s2}{fr\PYZus{}aliment}\PY{l+s+s2}{\PYZdq{}}\PY{p}{]}\PY{o}{.}\PY{n}{cumsum}\PY{p}{(}\PY{p}{)}

\PY{n}{freq\PYZus{}aliment}
\end{Verbatim}
\end{tcolorbox}

            \begin{tcolorbox}[breakable, size=fbox, boxrule=.5pt, pad at break*=1mm, opacityfill=0]
\prompt{Out}{outcolor}{24}{\boxspacing}
\begin{Verbatim}[commandchars=\\\{\}]
    Gastos com Alimentação  fi\_aliment  Fi\_aliment  fr\_aliment  Fr\_aliment
0                      1.5           1           1        0.04        0.04
1                      2.0           1           2        0.04        0.08
2                      6.0           1           3        0.04        0.12
3                      7.0           1           4        0.04        0.16
4                      8.0           1           5        0.04        0.20
5                     10.0           2           7        0.08        0.28
6                     12.0           1           8        0.04        0.32
7                     15.0           1           9        0.04        0.36
8                     20.0           2          11        0.08        0.44
9                     25.0           2          13        0.08        0.52
10                    30.0           2          15        0.08        0.60
11                    35.0           1          16        0.04        0.64
12                    40.0           5          21        0.20        0.84
13                    50.0           3          24        0.12        0.96
14                    60.0           1          25        0.04        1.00
\end{Verbatim}
\end{tcolorbox}
        
    \begin{tcolorbox}[breakable, size=fbox, boxrule=1pt, pad at break*=1mm,colback=cellbackground, colframe=cellborder]
\prompt{In}{incolor}{25}{\boxspacing}
\begin{Verbatim}[commandchars=\\\{\}]
\PY{n}{media\PYZus{}aliment} \PY{o}{=} \PY{p}{(}
    \PY{n}{freq\PYZus{}aliment}\PY{p}{[}\PY{l+s+s2}{\PYZdq{}}\PY{l+s+s2}{Gastos com Alimentação}\PY{l+s+s2}{\PYZdq{}}\PY{p}{]} \PY{o}{*} \PY{n}{freq\PYZus{}aliment}\PY{p}{[}\PY{l+s+s2}{\PYZdq{}}\PY{l+s+s2}{fr\PYZus{}aliment}\PY{l+s+s2}{\PYZdq{}}\PY{p}{]}
\PY{p}{)}\PY{o}{.}\PY{n}{sum}\PY{p}{(}\PY{p}{)}

\PY{n}{var\PYZus{}aliment} \PY{o}{=} \PY{p}{(}
    \PY{n}{freq\PYZus{}aliment}\PY{p}{[}\PY{l+s+s2}{\PYZdq{}}\PY{l+s+s2}{fi\PYZus{}aliment}\PY{l+s+s2}{\PYZdq{}}\PY{p}{]}
    \PY{o}{*} \PY{p}{(}\PY{n}{freq\PYZus{}aliment}\PY{p}{[}\PY{l+s+s2}{\PYZdq{}}\PY{l+s+s2}{Gastos com Alimentação}\PY{l+s+s2}{\PYZdq{}}\PY{p}{]} \PY{o}{\PYZhy{}} \PY{n}{media\PYZus{}aliment}\PY{p}{)}\PY{o}{*}\PY{o}{*}\PY{l+m+mi}{2}
\PY{p}{)}\PY{o}{.}\PY{n}{sum}\PY{p}{(}\PY{p}{)} \PY{o}{/} \PY{n}{n\PYZus{}aliment}

\PY{n}{std\PYZus{}aliment} \PY{o}{=} \PY{n}{np}\PY{o}{.}\PY{n}{sqrt}\PY{p}{(}\PY{n}{var\PYZus{}aliment}\PY{p}{)}

\PY{n}{cv\PYZus{}aliment} \PY{o}{=} \PY{n}{std\PYZus{}aliment} \PY{o}{/} \PY{n}{media\PYZus{}aliment} \PY{o}{*} \PY{l+m+mi}{100}

\PY{n}{g\PYZus{}aliment} \PY{o}{=} \PY{p}{(}
    \PY{n}{freq\PYZus{}aliment}\PY{p}{[}\PY{l+s+s2}{\PYZdq{}}\PY{l+s+s2}{fi\PYZus{}aliment}\PY{l+s+s2}{\PYZdq{}}\PY{p}{]}
    \PY{o}{*} \PY{p}{(}\PY{n}{freq\PYZus{}aliment}\PY{p}{[}\PY{l+s+s2}{\PYZdq{}}\PY{l+s+s2}{Gastos com Alimentação}\PY{l+s+s2}{\PYZdq{}}\PY{p}{]} \PY{o}{\PYZhy{}} \PY{n}{media\PYZus{}aliment}\PY{p}{)}\PY{o}{*}\PY{o}{*}\PY{l+m+mi}{3}
\PY{p}{)}\PY{o}{.}\PY{n}{sum}\PY{p}{(}\PY{p}{)} \PY{o}{/} \PY{p}{(}\PY{n}{n\PYZus{}aliment} \PY{o}{*} \PY{n}{std\PYZus{}aliment}\PY{o}{*}\PY{o}{*}\PY{l+m+mi}{3}\PY{p}{)}

\PY{n+nb}{print}\PY{p}{(}\PY{l+s+sa}{f}\PY{l+s+s1}{\PYZsq{}}\PY{l+s+s1}{O valor da média dos gastos com alimentação é: }\PY{l+s+si}{\PYZob{}}\PY{n}{media\PYZus{}aliment}\PY{l+s+si}{:}\PY{l+s+s1}{.2f}\PY{l+s+si}{\PYZcb{}}\PY{l+s+s1}{\PYZsq{}}\PY{p}{)}
\PY{n+nb}{print}\PY{p}{(}\PY{l+s+sa}{f}\PY{l+s+s1}{\PYZsq{}}\PY{l+s+s1}{O valor da variância dos gastos com alimentação é: }\PY{l+s+si}{\PYZob{}}\PY{n}{var\PYZus{}aliment}\PY{l+s+si}{:}\PY{l+s+s1}{.2f}\PY{l+s+si}{\PYZcb{}}\PY{l+s+s1}{\PYZsq{}}\PY{p}{)}
\PY{n+nb}{print}\PY{p}{(}\PY{l+s+sa}{f}\PY{l+s+s1}{\PYZsq{}}\PY{l+s+s1}{O valor do desvio padrão dos gastos com alimentação é: }\PY{l+s+si}{\PYZob{}}\PY{n}{std\PYZus{}aliment}\PY{l+s+si}{:}\PY{l+s+s1}{.2f}\PY{l+s+si}{\PYZcb{}}\PY{l+s+s1}{\PYZsq{}}\PY{p}{)}
\PY{n+nb}{print}\PY{p}{(}\PY{l+s+sa}{f}\PY{l+s+s1}{\PYZsq{}}\PY{l+s+s1}{O valor do coeficiente de variação é: }\PY{l+s+si}{\PYZob{}}\PY{n}{cv\PYZus{}aliment}\PY{l+s+si}{:}\PY{l+s+s1}{.2f}\PY{l+s+si}{\PYZcb{}}\PY{l+s+s1}{\PYZpc{}}\PY{l+s+s1}{\PYZsq{}}\PY{p}{)}
\PY{n+nb}{print}\PY{p}{(}\PY{l+s+sa}{f}\PY{l+s+s1}{\PYZsq{}}\PY{l+s+s1}{O valor da assimetria é: }\PY{l+s+si}{\PYZob{}}\PY{n}{g\PYZus{}aliment}\PY{l+s+si}{:}\PY{l+s+s1}{.2f}\PY{l+s+si}{\PYZcb{}}\PY{l+s+s1}{\PYZsq{}}\PY{p}{)}
\end{Verbatim}
\end{tcolorbox}

    \begin{Verbatim}[commandchars=\\\{\}]
O valor da média dos gastos com alimentação é: 26.66
O valor da variância dos gastos com alimentação é: 285.21
O valor do desvio padrão dos gastos com alimentação é: 16.89
O valor do coeficiente de variação é: 63.35\%
O valor da assimetria é: 0.16
    \end{Verbatim}

    \subsubsection{Quantis dos gastos com
Alimentação.}\label{quantis-dos-gastos-com-alimentauxe7uxe3o.}

    \begin{tcolorbox}[breakable, size=fbox, boxrule=1pt, pad at break*=1mm,colback=cellbackground, colframe=cellborder]
\prompt{In}{incolor}{26}{\boxspacing}
\begin{Verbatim}[commandchars=\\\{\}]
\PY{n}{aliment\PYZus{}ord} \PY{o}{=} \PY{n+nb}{sorted}\PY{p}{(}\PY{n}{aliment}\PY{p}{)}

\PY{n}{mediana\PYZus{}aliment} \PY{o}{=} \PY{n}{aliment\PYZus{}ord}\PY{p}{[}\PY{l+m+mi}{12}\PY{p}{]}
\PY{n}{q1\PYZus{}aliment} \PY{o}{=} \PY{p}{(}\PY{n}{aliment\PYZus{}ord}\PY{p}{[}\PY{l+m+mi}{5}\PY{p}{]} \PY{o}{+} \PY{n}{aliment\PYZus{}ord}\PY{p}{[}\PY{l+m+mi}{6}\PY{p}{]}\PY{p}{)} \PY{o}{/} \PY{l+m+mi}{2}
\PY{n}{q3\PYZus{}aliment} \PY{o}{=} \PY{p}{(}\PY{n}{aliment\PYZus{}ord}\PY{p}{[}\PY{l+m+mi}{18}\PY{p}{]} \PY{o}{+} \PY{n}{aliment\PYZus{}ord}\PY{p}{[}\PY{l+m+mi}{19}\PY{p}{]}\PY{p}{)} \PY{o}{/} \PY{l+m+mi}{2}

\PY{n+nb}{print}\PY{p}{(}\PY{n}{aliment\PYZus{}ord}\PY{p}{)}
\PY{n+nb}{print}\PY{p}{(}\PY{l+s+sa}{f}\PY{l+s+s1}{\PYZsq{}}\PY{l+s+s1}{A mediana dos gastos com alimentação é: }\PY{l+s+si}{\PYZob{}}\PY{n}{mediana\PYZus{}aliment}\PY{l+s+si}{:}\PY{l+s+s1}{.2f}\PY{l+s+si}{\PYZcb{}}\PY{l+s+s1}{\PYZsq{}}\PY{p}{)}
\PY{n+nb}{print}\PY{p}{(}\PY{l+s+sa}{f}\PY{l+s+s1}{\PYZsq{}}\PY{l+s+s1}{O valor de Q1 dos gastos com alimentação é: }\PY{l+s+si}{\PYZob{}}\PY{n}{q1\PYZus{}aliment}\PY{l+s+si}{:}\PY{l+s+s1}{.2f}\PY{l+s+si}{\PYZcb{}}\PY{l+s+s1}{\PYZsq{}}\PY{p}{)}
\PY{n+nb}{print}\PY{p}{(}\PY{l+s+sa}{f}\PY{l+s+s1}{\PYZsq{}}\PY{l+s+s1}{O valor de Q3 dos gastos com alimentação é: }\PY{l+s+si}{\PYZob{}}\PY{n}{q3\PYZus{}aliment}\PY{l+s+si}{:}\PY{l+s+s1}{.2f}\PY{l+s+si}{\PYZcb{}}\PY{l+s+s1}{\PYZsq{}}\PY{p}{)}
\end{Verbatim}
\end{tcolorbox}

    \begin{Verbatim}[commandchars=\\\{\}]
[1.5, 2, 6, 7, 8, 10, 10, 12, 15, 20, 20, 25, 25, 30, 30, 35, 40, 40, 40, 40,
40, 50, 50, 50, 60]
A mediana dos gastos com alimentação é: 25.00
O valor de Q1 dos gastos com alimentação é: 10.00
O valor de Q3 dos gastos com alimentação é: 40.00
    \end{Verbatim}

    \begin{tcolorbox}[breakable, size=fbox, boxrule=1pt, pad at break*=1mm,colback=cellbackground, colframe=cellborder]
\prompt{In}{incolor}{27}{\boxspacing}
\begin{Verbatim}[commandchars=\\\{\}]
\PY{n}{plt}\PY{o}{.}\PY{n}{boxplot}\PY{p}{(}\PY{n}{aliment}\PY{p}{,} \PY{n}{orientation}\PY{o}{=}\PY{l+s+s1}{\PYZsq{}}\PY{l+s+s1}{horizontal}\PY{l+s+s1}{\PYZsq{}}\PY{p}{,} \PY{n}{showmeans}\PY{o}{=}\PY{k+kc}{True}\PY{p}{)}
\PY{n}{plt}\PY{o}{.}\PY{n}{show}\PY{p}{(}\PY{p}{)}
\end{Verbatim}
\end{tcolorbox}

    \begin{center}
    \adjustimage{max size={0.9\linewidth}{0.9\paperheight}}{Exercicios_estatistica_rev_verificada_files/Exercicios_estatistica_rev_verificada_50_0.png}
    \end{center}
    { \hspace*{\fill} \\}
    
    8- Os valores da precipitação (em mm) registrada em uma Estação
Meteorológica, nos 31 dias de janeiro de 2014, estão listados na tabela:

{\def\LTcaptype{table} % do not increment counter
\begin{longtable}[]{@{}rr@{}}
\toprule\noalign{}
Dia & Precipitação \\
\midrule\noalign{}
\endhead
\bottomrule\noalign{}
\endlastfoot
1 & 0.0 \\
2 & 0.0 \\
3 & 0.0 \\
4 & 0.0 \\
5 & 4.7 \\
6 & 0.6 \\
7 & 17.2 \\
8 & 1.4 \\
9 & 11.2 \\
10 & 1.0 \\
11 & 3.8 \\
12 & 0.3 \\
13 & 0.0 \\
14 & 0.0 \\
15 & 0.5 \\
16 & 7.0 \\
17 & 0.0 \\
18 & 0.0 \\
19 & 3.3 \\
20 & 7.6 \\
21 & 0.9 \\
22 & 0.3 \\
23 & 18.2 \\
24 & 4.0 \\
25 & 4.6 \\
26 & 22.0 \\
27 & 15.6 \\
28 & 0.0 \\
29 & 3.4 \\
30 & 0.0 \\
31 & 0.0 \\
\end{longtable}
}

    \begin{itemize}
\tightlist
\item
  Escolha a melhor forma de representar os dados: representação
  individual ou por classes. Justifique.
\item
  Apresente a distribuição de frequências relativa e a distribuição de
  frequências relativa acumulada.
\item
  Apresente estatísticas de valores centrais.
\item
  Apresente estatísticas de dispersão.
\item
  Apresente estatísticas de forma. - Apresente o boxplot (gráfico de
  caixa) da variável aleatória.
\item
  Calcule os valores da variável aleatória correspondentes ao 1º, 2º e
  3º quartis.
\end{itemize}

    

    \begin{tcolorbox}[breakable, size=fbox, boxrule=1pt, pad at break*=1mm,colback=cellbackground, colframe=cellborder]
\prompt{In}{incolor}{28}{\boxspacing}
\begin{Verbatim}[commandchars=\\\{\}]
\PY{n}{dia} \PY{o}{=} \PY{n+nb}{list}\PY{p}{(}\PY{n+nb}{range}\PY{p}{(}\PY{l+m+mi}{1}\PY{p}{,} \PY{l+m+mi}{32}\PY{p}{)}\PY{p}{)}

\PY{n}{precipita} \PY{o}{=} \PY{p}{[}\PY{l+m+mf}{0.0}\PY{p}{,} \PY{l+m+mf}{0.0}\PY{p}{,} \PY{l+m+mf}{0.0}\PY{p}{,} \PY{l+m+mf}{0.0}\PY{p}{,} \PY{l+m+mf}{4.7}\PY{p}{,} \PY{l+m+mf}{0.6}\PY{p}{,} \PY{l+m+mf}{17.2}\PY{p}{,} \PY{l+m+mf}{1.4}\PY{p}{,} \PY{l+m+mf}{11.2}\PY{p}{,}
             \PY{l+m+mf}{1.0}\PY{p}{,} \PY{l+m+mf}{3.8}\PY{p}{,} \PY{l+m+mf}{0.3}\PY{p}{,} \PY{l+m+mf}{0.0}\PY{p}{,} \PY{l+m+mf}{0.0}\PY{p}{,} \PY{l+m+mf}{0.5}\PY{p}{,} \PY{l+m+mi}{7}\PY{p}{,} \PY{l+m+mf}{0.0}\PY{p}{,} \PY{l+m+mf}{0.0}\PY{p}{,} \PY{l+m+mf}{3.3}\PY{p}{,} \PY{l+m+mf}{7.6}\PY{p}{,} \PY{l+m+mf}{0.9}\PY{p}{,}
             \PY{l+m+mf}{0.3}\PY{p}{,} \PY{l+m+mf}{18.2}\PY{p}{,} \PY{l+m+mi}{4}\PY{p}{,} \PY{l+m+mf}{4.6}\PY{p}{,} \PY{l+m+mi}{22}\PY{p}{,} \PY{l+m+mf}{15.6}\PY{p}{,} \PY{l+m+mi}{0}\PY{p}{,} \PY{l+m+mf}{3.4}\PY{p}{,} \PY{l+m+mf}{0.0}\PY{p}{,} \PY{l+m+mf}{0.0}\PY{p}{]}
\PY{n}{n\PYZus{}precipita} \PY{o}{=} \PY{n+nb}{len}\PY{p}{(}\PY{n}{precipita}\PY{p}{)}

\PY{n}{dados\PYZus{}precipita} \PY{o}{=} \PY{n}{pd}\PY{o}{.}\PY{n}{DataFrame}\PY{p}{(}\PY{p}{\PYZob{}}
    \PY{l+s+s2}{\PYZdq{}}\PY{l+s+s2}{Dia}\PY{l+s+s2}{\PYZdq{}}\PY{p}{:} \PY{n}{dia}\PY{p}{,} \PY{l+s+s2}{\PYZdq{}}\PY{l+s+s2}{Precipitação}\PY{l+s+s2}{\PYZdq{}}\PY{p}{:} \PY{n}{precipita}
\PY{p}{\PYZcb{}}\PY{p}{)}


\PY{n}{freq\PYZus{}precipita} \PY{o}{=} \PY{p}{(}
    \PY{n}{dados\PYZus{}precipita}\PY{p}{[}\PY{l+s+s2}{\PYZdq{}}\PY{l+s+s2}{Precipitação}\PY{l+s+s2}{\PYZdq{}}\PY{p}{]}
    \PY{o}{.}\PY{n}{value\PYZus{}counts}\PY{p}{(}\PY{p}{)}
    \PY{o}{.}\PY{n}{sort\PYZus{}index}\PY{p}{(}\PY{p}{)}
    \PY{o}{.}\PY{n}{reset\PYZus{}index}\PY{p}{(}\PY{p}{)}
\PY{p}{)}

\PY{n}{freq\PYZus{}precipita}\PY{o}{.}\PY{n}{columns} \PY{o}{=} \PY{p}{[}\PY{l+s+s2}{\PYZdq{}}\PY{l+s+s2}{Precipitação}\PY{l+s+s2}{\PYZdq{}}\PY{p}{,} \PY{l+s+s2}{\PYZdq{}}\PY{l+s+s2}{fi\PYZus{}precipita}\PY{l+s+s2}{\PYZdq{}}\PY{p}{]}

\PY{n}{freq\PYZus{}precipita}\PY{p}{[}\PY{l+s+s2}{\PYZdq{}}\PY{l+s+s2}{Fi\PYZus{}precipita}\PY{l+s+s2}{\PYZdq{}}\PY{p}{]} \PY{o}{=} \PY{n}{freq\PYZus{}precipita}\PY{p}{[}\PY{l+s+s2}{\PYZdq{}}\PY{l+s+s2}{fi\PYZus{}precipita}\PY{l+s+s2}{\PYZdq{}}\PY{p}{]}\PY{o}{.}\PY{n}{cumsum}\PY{p}{(}\PY{p}{)}
\PY{n}{freq\PYZus{}precipita}\PY{p}{[}\PY{l+s+s2}{\PYZdq{}}\PY{l+s+s2}{fr\PYZus{}precipita}\PY{l+s+s2}{\PYZdq{}}\PY{p}{]} \PY{o}{=} \PY{n}{freq\PYZus{}precipita}\PY{p}{[}\PY{l+s+s2}{\PYZdq{}}\PY{l+s+s2}{fi\PYZus{}precipita}\PY{l+s+s2}{\PYZdq{}}\PY{p}{]} \PY{o}{/} \PY{n}{n\PYZus{}precipita}
\PY{n}{freq\PYZus{}precipita}\PY{p}{[}\PY{l+s+s2}{\PYZdq{}}\PY{l+s+s2}{Fr\PYZus{}precipita}\PY{l+s+s2}{\PYZdq{}}\PY{p}{]} \PY{o}{=} \PY{n}{freq\PYZus{}precipita}\PY{p}{[}\PY{l+s+s2}{\PYZdq{}}\PY{l+s+s2}{fr\PYZus{}precipita}\PY{l+s+s2}{\PYZdq{}}\PY{p}{]}\PY{o}{.}\PY{n}{cumsum}\PY{p}{(}\PY{p}{)}
\end{Verbatim}
\end{tcolorbox}

    \begin{tcolorbox}[breakable, size=fbox, boxrule=1pt, pad at break*=1mm,colback=cellbackground, colframe=cellborder]
\prompt{In}{incolor}{29}{\boxspacing}
\begin{Verbatim}[commandchars=\\\{\}]
\PY{c+c1}{\PYZsh{} Média}
\PY{n}{media\PYZus{}precipita} \PY{o}{=} \PY{p}{(}
    \PY{n}{freq\PYZus{}precipita}\PY{p}{[}\PY{l+s+s2}{\PYZdq{}}\PY{l+s+s2}{fr\PYZus{}precipita}\PY{l+s+s2}{\PYZdq{}}\PY{p}{]}
    \PY{o}{*} \PY{n}{freq\PYZus{}precipita}\PY{p}{[}\PY{l+s+s2}{\PYZdq{}}\PY{l+s+s2}{Precipitação}\PY{l+s+s2}{\PYZdq{}}\PY{p}{]}
\PY{p}{)}\PY{o}{.}\PY{n}{sum}\PY{p}{(}\PY{p}{)}

\PY{c+c1}{\PYZsh{} Variância populacional pela tabela de frequências}
\PY{n}{var\PYZus{}precipita} \PY{o}{=} \PY{p}{(}
    \PY{n}{freq\PYZus{}precipita}\PY{p}{[}\PY{l+s+s2}{\PYZdq{}}\PY{l+s+s2}{fi\PYZus{}precipita}\PY{l+s+s2}{\PYZdq{}}\PY{p}{]}
    \PY{o}{*} \PY{p}{(}\PY{n}{freq\PYZus{}precipita}\PY{p}{[}\PY{l+s+s2}{\PYZdq{}}\PY{l+s+s2}{Precipitação}\PY{l+s+s2}{\PYZdq{}}\PY{p}{]} \PY{o}{\PYZhy{}} \PY{n}{media\PYZus{}precipita}\PY{p}{)}\PY{o}{*}\PY{o}{*}\PY{l+m+mi}{2}
\PY{p}{)}\PY{o}{.}\PY{n}{sum}\PY{p}{(}\PY{p}{)} \PY{o}{/} \PY{n}{n\PYZus{}precipita}

\PY{n}{moda\PYZus{}precipita} \PY{o}{=} \PY{n}{freq\PYZus{}precipita}\PY{o}{.}\PY{n}{loc}\PY{p}{[}
    \PY{n}{freq\PYZus{}precipita}\PY{p}{[}\PY{l+s+s2}{\PYZdq{}}\PY{l+s+s2}{fi\PYZus{}precipita}\PY{l+s+s2}{\PYZdq{}}\PY{p}{]}\PY{o}{.}\PY{n}{idxmax}\PY{p}{(}\PY{p}{)}\PY{p}{,}
    \PY{l+s+s2}{\PYZdq{}}\PY{l+s+s2}{Precipitação}\PY{l+s+s2}{\PYZdq{}}
\PY{p}{]}

\PY{n}{precipita\PYZus{}ord} \PY{o}{=} \PY{n+nb}{sorted}\PY{p}{(}\PY{n}{precipita}\PY{p}{)}
\PY{n}{mediana\PYZus{}precipita} \PY{o}{=} \PY{n}{precipita\PYZus{}ord}\PY{p}{[}\PY{l+m+mi}{15}\PY{p}{]}

\PY{n+nb}{print}\PY{p}{(}\PY{l+s+sa}{f}\PY{l+s+s1}{\PYZsq{}}\PY{l+s+s1}{A média da precipitação é: }\PY{l+s+si}{\PYZob{}}\PY{n}{media\PYZus{}precipita}\PY{l+s+si}{:}\PY{l+s+s1}{.2f}\PY{l+s+si}{\PYZcb{}}\PY{l+s+s1}{\PYZsq{}}\PY{p}{)}
\PY{n+nb}{print}\PY{p}{(}\PY{l+s+sa}{f}\PY{l+s+s1}{\PYZsq{}}\PY{l+s+s1}{A variância da precipitação é: }\PY{l+s+si}{\PYZob{}}\PY{n}{var\PYZus{}precipita}\PY{l+s+si}{:}\PY{l+s+s1}{.2f}\PY{l+s+si}{\PYZcb{}}\PY{l+s+s1}{\PYZsq{}}\PY{p}{)}
\PY{n+nb}{print}\PY{p}{(}\PY{l+s+sa}{f}\PY{l+s+s1}{\PYZsq{}}\PY{l+s+s1}{A moda da precipitação é: }\PY{l+s+si}{\PYZob{}}\PY{n}{moda\PYZus{}precipita}\PY{l+s+si}{:}\PY{l+s+s1}{.2f}\PY{l+s+si}{\PYZcb{}}\PY{l+s+s1}{\PYZsq{}}\PY{p}{)}
\PY{n+nb}{print}\PY{p}{(}\PY{l+s+sa}{f}\PY{l+s+s1}{\PYZsq{}}\PY{l+s+s1}{A mediana da precipitação é: }\PY{l+s+si}{\PYZob{}}\PY{n}{mediana\PYZus{}precipita}\PY{l+s+si}{:}\PY{l+s+s1}{.2f}\PY{l+s+si}{\PYZcb{}}\PY{l+s+s1}{\PYZsq{}}\PY{p}{)}
\end{Verbatim}
\end{tcolorbox}

    \begin{Verbatim}[commandchars=\\\{\}]
A média da precipitação é: 4.12
A variância da precipitação é: 37.49
A moda da precipitação é: 0.00
A mediana da precipitação é: 0.90
    \end{Verbatim}

    \begin{tcolorbox}[breakable, size=fbox, boxrule=1pt, pad at break*=1mm,colback=cellbackground, colframe=cellborder]
\prompt{In}{incolor}{ }{\boxspacing}
\begin{Verbatim}[commandchars=\\\{\}]

\end{Verbatim}
\end{tcolorbox}


    % Add a bibliography block to the postdoc
    
    
    
\end{document}

